\documentclass[10pt]{amsart}
\usepackage[a4paper,margin=23mm]{geometry}
\usepackage{amsmath,amssymb,amsthm,mathtools,hyperref,graphicx,xspace,letltxmacro,xpatch}
\usepackage[english]{babel}
\newtheorem{theo}{Theorem}[section]
\newtheorem{lemm}[theo]{Lemma}\newtheorem{prop}[theo]{Proposition}\newtheorem{coro}[theo]{Corollary}\newtheorem{conj}[theo]{Conjecture}
\theoremstyle{definition}\newtheorem{defi}[theo]{Definition}\newtheorem{rema}[theo]{Remark}\newtheorem{exam}[theo]{Example}\newtheorem{exem}[theo]{Example}\newtheorem{nota}[theo]{Notation}
\newcommand{\ie}{i.e.}\newcommand{\eg}{e.g.}\newcommand{\resp}{resp.}\newcommand{\skpt}{\par\smallskip\noindent}
\emergencystretch=3em\pagestyle{plain}





%\AtBeginDocument{
%\newtheorem{condition}[cdrthm]{Condition} }
\newcommand{\clI}{\mathcal{I}}
\newcommand{\clC}{\mathcal{C}}
\newcommand{\clL}{\mathcal{L}}
\newcommand{\clN}{\mathcal{N}}
\newcommand{\clH}{\mathcal{H}}
\newcommand{\clF}{\mathcal{F}}
\newcommand{\clB}{\mathcal{B}}
\newcommand{\bbS}{\mathbb{S}}
\newcommand{\bbR}{\mathbb{R}}
\newcommand{\bbP}{\mathbb{P}}
\newcommand{\bbE}{\mathbb{E}}
\newcommand{\bbN}{\mathbb{N}}
\newcommand{\fkF}{\mathfrak{F}}
\newcommand{\fkB}{\mathfrak{B}}

\newcommand{\rVar}{\mathrm{Var}}
\DeclareMathOperator{\Var}{Var}
%%%%%%%%%%%%%%%%%%%%%%%%%%%%%%%%%%%%%%%%%%%%%%%%%%%%%%%%%%%%%%%%%%%%%%%%%%%%%%%%%%%%%%%%%%%%%%%%%%%%%%%%%%%
\graphicspath{{./figures/}}

\newcommand*{\mk}{\mkern -1mu}
\newcommand*{\Mk}{\mkern -2mu}
\newcommand*{\mK}{\mkern 1mu}
\newcommand*{\MK}{\mkern 2mu}

\hypersetup{urlcolor=purple, linkcolor=blue, citecolor=red}

\newcommand*{\romanenumi}{\renewcommand*{\theenumi}{\roman{enumi}}}
\newcommand*{\Romanenumi}{\renewcommand*{\theenumi}{\Roman{enumi}}}
\newcommand*{\alphenumi}{\renewcommand*{\theenumi}{\alph{enumi}}}
\newcommand*{\Alphenumi}{\renewcommand*{\theenumi}{\Alph{enumi}}}
\let\oldtilde\tilde
\renewcommand*{\tilde}[1]{\mathchoice{\widetilde{#1}}{\widetilde{#1}}{\oldtilde{#1}}{\oldtilde{#1}}}
\let\oldhat\hat
\renewcommand*{\hat}[1]{\mathchoice{\widehat{#1}}{\widehat{#1}}{\oldhat{#1}}{\oldhat{#1}}}
%%%%%%%%%%%%%%%%%%%%%%%%%%%%%%%%%%%%%%%%%%%%%%%%%%%%%%%%%%%%%%%%%%%%%%%%%%%%%%%%%%%%%%%%%%%%%%%%%%%%%%%%%%%%%%%%%%%%%


\NewDocumentEnvironment{enonce}{ O{} m O{} }{\par\smallskip\refstepcounter{theo}\noindent\textbf{#2~\thetheo\if\relax\detokenize{#3}\relax\else\ (#3)\fi.}\ }{\par\smallskip}
\title{Long-memory noncentral limits for level-curve lengths of time-dependent spherical Gaussian fields}
\author{Marinucci, Domenico, Rossi, Maurizia, Vidotto, Anna}
\date{doi:10.5802/ahl.206; CC BY 4.0}
\begin{document}
\maketitle
\noindent\textit{Editorial note.} Only selected theorem counted; complete author mathematical body and original reference keys retained. Source-local supporting results and appendices remain uncounted context.\par\medskip
\noindent\textit{Allowed standard prerequisites (editorial).}\begin{itemize}
\item Standard functional analysis, measure theory and the author-cited classical theorems. Standard prior results stay cited; complete source-local proofs and definitions retained.
\end{itemize}\medskip

\section{Background and notation}

The analysis of level curves for random fields is a very classical topic in stochastic geometry. In particular, many efforts have focussed on the investigation of level-zero curves (i.e., nodal lines) in the case of random eigenfunctions, in the high-frequency regime where eigenvalues are assumed to diverge to infinity; see for instance~\cite{MPRW16, NPR19, PV20,Wig}, or more generally~\cite{Wig22} for a recent overview. In the same high-energy regime, other functionals for random eigenfunctions (including excursion area, the Euler--Poincaré characteristic, the number of critical points) have also been widely investigated, see for instance~\cite {M2022}; on the other hand, these same functionals have also been considered by different authors in the asymptotic regime where the spatial domain of the field is assumed as growing, notable examples being~\cite{KL01} (for level curves) and~\cite{EL2016} (for the Euler--Poincaré characteristic).

Our purpose in this paper is to study the behaviour of level curves under a different asymptotic regime than so far considered, namely for sphere cross-time random fields and taking into account the averaged fluctuations over time around the expected value (a similar framework was considered for the case of the excursion area in~\cite{MRV:20}). Our asymptotic results share some analogies with the different settings that we mentioned above, but they also show very important differences that we shall discuss below in greater detail. While this paper only focusses on theoretical aspects, it is really not difficult to envisage application areas where sphere cross-time random fields emerge very naturally, some examples being atmospheric and climate data (the sphere representing the surface of the Earth, see~\cite {Ch:17}).

\subsection{Time-dependent spherical random fields}

We start by recalling the notion of space-time spherical random field, along with the corresponding spectral representation, which allows the characterization of long and short range dependence properties. Our assumptions and discussion is close to the one that can be found in~\cite {MRV:20}.

More precisely, let us take a probability space $(\Omega,\fkF,
\bbP)$ and denote by $\bbE$ the expectation under $\bbP$: all random objects in this manuscript are defined on this common probability space, unless otherwise specified. Let $\bbS^{2}$ denote the two-dimensional unit sphere with the round metric, usually written in the form
\begin{equation*}
ds^{2}=d\theta ^{2}+\sin ^{2}\theta d\varphi ^{2}
\end{equation*}
for standard spherical coordinates $(\theta,\varphi)$, where $\theta $ is the colatitude and $\varphi $ the longitude. We will denote by $\Delta_{\bbS^2}$ the spherical Laplacian, in coordinates
\begin{equation*}
\Delta_{\bbS^2} = \frac{1}{\sin \theta} \frac{\partial}{\partial \theta} \left (\sin
\theta \frac{\partial}{\partial \theta} \right) + \frac{1}{\sin^2\theta}
\frac{\partial^2}{\partial \varphi^2}.
\end{equation*}
A space-time real-valued spherical random field $Z$ is a collection of real random variables indexed by $\bbS^{2}\times \bbR$
\begin{equation}\label{eq1}
Z=\left\{Z(x,t),\,x\in \bbS^{2},\,t\in \bbR\right\}
\end{equation}
such that the function $Z:\Omega \times \bbS^{2}\times \bbR
\rightarrow \bbR$ is $\fkF\otimes \fkB(\bbS ^{2}\times \bbR)$-measurable, $\fkB(\bbS^{2}\times\bbR)$ being the Borel $\sigma $-field of $\bbS^{2}\times
\bbR$. The following condition is standard.

\begin{enonce}{Condition}\label{basic}
The space-time real-valued spherical random field $Z$ in~\eqref{eq1} is
\begin{itemize}
\item Gaussian, i.e.~its finite dimensional distributions are Gaussian;

\item centered, that is, $\bbE[Z(x,t)] = 0$ for every $x\in \bbS ^2$, $t\in \bbR$;

\item isotropic (in space) and stationary (in time), namely
\begin{equation}\label{Gamma}
\bbE[Z(x,t) Z(y,s)] = \Gamma(\langle x,y\rangle, t-s)
\end{equation}
for every $x,y\in \bbS^2$, $t,s \in \bbR$, where $\langle x, y\rangle$ denotes the standard inner product in $\bbR^3$ and $
\Gamma:[-1,1]\times \bbR\to \bbR$ is a positive semidefinite function;

\item mean square continuous, i.e. $\Gamma$ is continuous.
\end{itemize}
\end{enonce}

Assumption~\ref{basic} collects the common background with basically all the previous literature, starting from~\cite{adlertaylor}, on the geometry of excursion sets of (time-varying) random fields on manifolds, see also~\cite {BP17, LO13, MRV:20}. \emph{From now on we assume that $Z$ in~\eqref{eq1} satisfies Assumption~\ref{basic}.}

Under Assumption~\ref{basic} it is well known (see e.g.~\cite[Theorem~3.3] {BP17} and~\cite[Theorem~3]{MM18}) that the covariance function $\Gamma$ in~\eqref{Gamma} of $Z$ in~\eqref{eq1} can be written as a uniformly convergent series of the form
\begin{equation}\label{eq2}
\Gamma(\eta, \tau) = \sum_{\ell=0}^{+\infty} \frac{2\ell+1}{4\pi} C_\ell(\tau) P_\ell(\eta)\,,\quad (\eta, \tau) \in [-1,1] \times \bbR
\,,
\end{equation}
where $\{ C_\ell, \ell\ge 0\}$ is a sequence of continuous positive semidefinite functions on the real line and $\{ P_\ell,
\ell\ge 0\}$ stands for the sequence of Legendre polynomials: $
\int_{-1}^1 P_{\ell}(t) P_{\ell^{\prime }}(t)\,dt= \frac{2}{2\ell+1}
\delta_{\ell}^{\ell^{\prime }}, $ $\delta_{\ell}^{\ell^{\prime }}$ denoting the Kronecker delta, see~\cite[Section~4.7]{Sze75}. Note that the uniform convergence of the series~\eqref{eq2} is equivalent to
\begin{equation*}
\sum_{\ell = 0}^{+\infty} \frac{2\ell+1}{4\pi} C_\ell(0) < +\infty
\end{equation*}
($C_\ell(0)\ge 0$ for every $\ell\ge 0$). \\
In addition, we assume from now that there exists the spatial gradient of $Z$, i.e., $\nabla Z = \{ \nabla Z(x,t), x\in \bbS^2, t\in \bbR\}$ is a measurable map on $\Omega\times \bbS^2\times \bbR$. We also require $\nabla Z$ to be mean square continuous. Note that $
\nabla Z $ is a centered Gaussian random field indexed by $\bbS^2\times \bbR $ whose covariance kernel is the spatial Hessian of $\Gamma$ in~\eqref{Gamma}.

\begin{enonce}{Condition}\label{regularT}
A.s., the random field $\nabla Z(\cdot,t)$ is in $\clC^0(\bbS^2)$ for every $t\in \bbR$.
\end{enonce}
\emph{From now on we assume that $Z$ in~\eqref{eq1} satisfies also Assumption~\ref{regularT}. } In particular, $\sum_{\ell=0}^{+\infty} (2\ell+1) C_\ell(0) \cdot \ell^2 < +\infty$.

Moreover, we assume that there exists the spherical Laplacian field $\Delta_{\bbS^2} Z$, i.e., $\Delta_{\bbS^2} Z = \{ \Delta_{\bbS^2} Z(x,t), x\in \bbS^2, t\in \bbR\}$ is a measurable map on $\Omega\times \bbS^2\times \bbR$. We also require $\Delta_{\bbS^2} Z$ to be mean square continuous. Hence,
\begin{equation}\label{eqConv}
\sum_{\ell = 0}^{+\infty} \frac{2\ell+1}{4\pi} C_\ell(0) \cdot \ell^4 < +\infty.
\end{equation}

In particular, for $k=1,2,3,4$ the following series involving the $k^{\rm th}$ derivative of Legendre polynomials uniformly converges
\begin{equation}\label{eqDer}
\frac{\partial^k}{\partial \eta^k}\Gamma(\eta, \tau)= \sum_{\ell=0}^{+\infty}
\frac{2\ell+1}{4\pi} C_\ell(\tau) \frac{\partial^k}{\partial \eta^k} P_\ell(\eta),\quad (\eta, \tau) \in [-1,1] \times \bbR.
\end{equation}
(In other words, the spatial derivatives up to order four of the covariance function admit a series representation of the form~\eqref{eq2} with the Legendre polynomials replaced by their derivatives.)



Let us now introduce some more notation: first denote by $\{ Y_{\ell,m}, \ell \ge 0, m=-\ell,\, \dots,\, \ell\}$ the standard real orthonormal basis of spherical harmonics~\cite[Section~3.4]{MaPeCUP} for $ L^2(\bbS^2)$, then define for $\ell\in \bbN, m=-\ell,\,\dots,\,\ell$,
\begin{equation}\label{eq3}
a_{\ell,m}(t) := \int_{\bbS^2} Z(x,t) Y_{\ell,m}(x)\,dx,\quad t\in
\bbR.
\end{equation}
From~\eqref{eq3} we deduce that $\{ a_{\ell,m}, \ell \ge 0, m=-\ell,\,
\dots,\,\ell\}$ is a family of independent, stationary, centered, Gaussian processes on the real line such that for every $t,s \in \bbR$
\begin{equation*}
\bbE[a_{\ell,m}(t) a_{\ell,m}(s)] = C_\ell(t-s).
\end{equation*}
The spectral representation~\eqref{eq2} for $\Gamma$ allows to deduce the so-called Karhunen--Lo\`eve expansion for $Z$:
\begin{equation}\label{eqSeries}
Z(x,t) = \sum_{\ell =0}^{+\infty} \sum_{m=-\ell}^\ell a_{\ell,m}(t) Y_{\ell,m}(x),
\end{equation}
where the stochastic processes $a_{\ell,m}$ are defined as in~\eqref{eq3}, and the series~\eqref{eqSeries} converges in $L^2(\Omega \times \bbS^2
\times [0,T])$ for any $T>0$. For $\bbS^2\ni x=(\theta_x, \varphi_x)$ we use the notation

\begin{equation*}
\partial_{1;x} = \frac{\partial}{\partial \theta}\bigg|_{\theta=\theta_x},\quad
\partial_{2;x}= \frac{1}{\sin\theta}\frac{\partial}{\partial}\bigg|_{\theta=\theta_x, \varphi=\varphi_x}.
\end{equation*}
Analogously, \eqref{eqDer} ensures that, for $j=1,2$,
\begin{equation}\label{eqSeries_gradient}
\partial_{j;x} Z(x,t) = \sum_{\ell =0}^{+\infty}\;\sum_{m=-\ell}^\ell a_{\ell,m}(t) \partial_{j;x} Y_{\ell,m}(x),
\end{equation}
where the convergence still holds in $L^2(\Omega \times \bbS^2 \times [0,T])$. Analogous series expansions hold in $L^2(\Omega \times \bbS^2 \times [0,T])$ for the spherical Laplacian $\partial^2_{1;x} Z + \partial^2_{2;x}Z$.

From now on we can restrict ourselves to
\begin{equation*}
\widetilde{\bbN} := \{ \ell \ge 0 : C_\ell(0)\ne 0\}
\end{equation*}
without loss of generality.

\subsubsection{Time-dependent random spherical eigenfunctions}
\label{sec-rsh}

Let us define
\begin{equation}\label{eqZell}
Z_\ell(x,t) := \sum_{m=-\ell}^\ell a_{\ell,m}(t) Y_{\ell,m}(x)\,, \quad (x,t)\in \bbS^2\times \bbR;
\end{equation}
by construction, $\{ Z_\ell, \ell\in \widetilde{\bbN}\}$ is a sequence of independent random fields and each $Z_\ell(\cdot, t)$ almost surely solves the Helmholtz equation $\Delta_{\bbS^2} Z_{\ell }(\cdot, t)+\lambda_\ell Z_{\ell }(\cdot, t)=0, $ where $\lambda_\ell:=\ell (\ell +1)$ is the $\ell^{\rm th}$ eigenvalue. For the sake of notational simplicity we will assume that (cf.~\eqref{eqConv})
\begin{equation}\label{somma1}
\sigma _{0}^{2}:=\bbE\left[Z^{2}(x,t)\right] =\sum_{\ell\,\in\,\widetilde{\bbN}} \bbE\left[Z_\ell(x,t)^2\right]=\sum_{\ell\,\in\,\widetilde{
\bbN}}\frac{2\ell +1}{4\pi }C_{\ell }(0)=1.
\end{equation}

\paragraph{Some conventions} 
From now on, $c\in (0, +\infty)$ will stand for a universal constant which may change from line to line. Let $\{ a_n, n\ge 0\}$, $\{ b_n, n\ge 0\}$ be two sequences of positive numbers: we will write $a_n \sim b_n$ if $a_n/b_n \to 1$ as $n\to +\infty$, $ a_n \approx b_n$ whenever $a_n/b_n \to c$, $a_n = o(b_n)$ if $a_n/b_n \to 0$, and finally $a_n = O(b_n)$ or equivalently $a_n \ll b_n$ if eventually $ a_n/b_n \le c$.


\subsection{Time dependence properties}

As in~\cite{MRV:20}, let us now define the family of symmetric real-valued functions $\{ g_\beta, \beta\in (0,1]\}$ as follows: for $\tau\in \bbR$
\begin{equation}\label{gbeta}
g_\beta(\tau) =
\begin{cases}
(1 + |\tau|)^{-\beta}\, & \text{if }\,\, \beta\in (0,1) \\
(1+|\tau|)^{-\alpha}\, & \text{if }\,\, \beta = 1
\end{cases}
\,,
\end{equation}
for some $\alpha \in [2,+\infty)$. We do believe that the assumption $\alpha
\in [2,+\infty)$ is not essential for the validity of our main findings; indeed it seems likely that it can be replaced with $\alpha \in (1, +\infty)$. Nevertheless, the current formulation is instrumental to take advantage of some technical results provided in~\cite{MRV:20}.

\begin{enonce}{Condition}\label{basic3}
There exists a sequence of numbers $\{ \beta_\ell\in (0,1], \ell
\in \widetilde{\bbN}\}$ and a sequence of real valued functions $\{ G_\ell, \ell
\in \widetilde{\bbN}\}$ such that
\begin{equation*}
C_\ell(\tau) = G_\ell(\tau) \cdot g_{\beta_\ell}(\tau), \qquad \ell\in
\widetilde{\bbN},
\end{equation*}
where $g_{\beta_\ell}$ is as in~\eqref{gbeta} and
\begin{equation*}
\sup_{\ell\,\in\,\widetilde{\bbN}} \left|\frac{G_\ell(\tau)}{G_\ell(0)} - 1\right| = o(1), \quad \text{as }\tau\to +\infty\,.
\end{equation*}
Moreover $0\in \widetilde{\bbN}$ (that is, $C_0(0)\ne 0$) and if $
\beta_0=1$ then
\begin{equation*}
\int_{\bbR} C_0 (\tau)\, d\tau > 0\,.
\end{equation*}
\end{enonce}

\emph{From now on we assume that Assumption~\ref{basic3} holds for the sequence $\{ C_\ell, \ell\in \widetilde{\bbN}\}$.} Note that $G_\ell(0)=C_\ell(0)$ for every $\ell\in \widetilde{\bbN}$.

Let $\ell\in \widetilde{\bbN}$. As discussed in~\cite{MRV:20}, the coefficient $\beta_{\ell }$ in Assumption~\ref{basic3} governs the \emph{ memory} of our processes; indeed, for $\beta_{\ell }=1$ (resp.~$
\beta_\ell\in (0,1)$) the covariance function $C_{\ell }$ is integrable on $\bbR$ (resp.~$\int_{\bbR} |C_\ell(\tau)|\,d\tau = +\infty$) and the corresponding process has so-called short (resp.~long) memory behavior (note that $C_\ell(0)$ is always non-negative but $C_\ell(\tau)$ need not be, for $\tau>0 $).

Clearly one could choose alternative parametrizations for $g_\beta(\tau)$, such as for instance
\begin{equation*}
g_\beta(\tau) =\left(1 + |\tau|^2\right)^{-\beta/2},\quad \text{or} \quad g_\beta(\tau) =\left(1 + |\tau|^{\gamma}\right)^{-\beta};
\end{equation*}
however, these choices obviously cannot change our results, as our condition is basically requiring that, for all $\ell$,
\begin{equation*}
\lim_{\tau\,\rightarrow\,\infty}\frac{C_\ell(\tau)}{C_\ell(0)\tau^{-\beta_\ell}} =1.
\end{equation*}
A possible generalizations would be to allow for the possibility of slowly-varying (at infinity) factors, i.e. to allow for autocorrelations of the form $ L(|\tau|)\tau^{-\beta}$, where $L(\cdot)$ is such that $\lim_{\tau\,\rightarrow\,\infty} L(|\tau|)/L(a|\tau|)=1$ for all $a>0$. These generalizations are common in the long memory literature but would not alter by any means the substance of our results, so we avoid to consider them for brevity's sake.



\section{Main Results}\label{sec-MR}

In this Section we introduce the problem and describe our main results. We start with some technical lemmas.

\subsection{The average boundary length}
\label{subsec-object}

Let $u\in \bbR$ be a threshold fixed from now on, for $t\in \bbR$ we consider the level set
\begin{equation*}
Z(\cdot, t)^{-1}(u):=\left\{ x\in \bbS^2 : Z(x,t) = u\right\}
\end{equation*}
which is a.s. a $\clC^1$ manifold of dimension~$1$ thanks to Ylvisaker's Lemma~\cite[Theorem~1.21]{AW:09} (see also Bulinskaya's Lemma~\cite[Proposition~1.20]{AW:09}) - note that here \emph{``a.s.'' may depend on $t$}. Indeed, for every $x\in Z(\cdot, t)^{-1}(u)$ the random vector $(Z(x,t), \nabla Z(x,t))$ is non-degenerate hence, except for a negligible subset of $\Omega$, say $\Omega_t^c (\bbP(\Omega_t^c)=0)$, the value $u$ is regular for $Z(\cdot, t)$, i.e. $\nabla Z(x, t)\ne 0$ for every $x$ such that $Z(x,t)=u$.

Let $t\in \bbR$ be fixed, on $\Omega_t (\bbP(\Omega_t)=1)$ we define
\begin{equation}\label{length}
\clL_u(t) := \clH^1\left(Z(\cdot, t)^{-1}(u)\right),
\end{equation}
i.e., the $1$-dimensional Hausdorff measure of the level set $Z(\cdot, t)^{-1}(u)$, that we refer to as the \emph{length of $u$-level curves at time $t$}.

By (time) stationarity of $Z$, the law of $\clL_u(t)$ does not depend on $t$, in particular $\bbE[\clL_u(t)]$ does not depend on $t$, and can be computed via the Kac-Rice formula~\cite[Theorem~6.8]{AW:09} or the Gaussian Kinematic Formula~\cite[Theorem~13.2.1]{adlertaylor} to be
\begin{equation}\label{meanGKF}
\bbE[\clL_u(t)] = \sigma_1 \cdot 2\pi e^{-u^2/2} = \bbE[\clL_u(0)],
\end{equation}
where for any $j=1,2$
\begin{equation}\label{sigma1}
\sigma _{1}^{2}:=\bbE\left[\partial _{j;x}Z(x,t)\,\partial _{j;x}Z(x,t)\right] =\sum_{\ell =0}^{\infty}\frac{2\ell +1}{4\pi }C_{\ell }(0)\frac{\ell (\ell +1)}{2},
\end{equation}
cf.~\eqref{eqConv}. See Lemma~\ref{lemcov} in Appendix~\ref{app-cov} for details on the covariance structure of the field $(Z, \nabla Z)$.

\begin{lemm}\label{lem_basic}
There exists $\tilde \Omega \subseteq \Omega$ such that $
\bbP(\tilde \Omega)=1$ and for every $\omega\in \tilde \Omega$ there exists $I(\omega)\subseteq [0,+\infty)$ whose complement is negligible ($\bbP(\tilde{\Omega}^c)=0$) such that the value $u$ is regular for $Z(\cdot, t)(\omega)$ for every $t\in I(\omega)$.
\end{lemm}

In view of Lemma~\ref{lem_basic}, whose proof is postponed to the Appendix~\ref{app-square}, we can define the family $\{ \clC_T(u), T>0\}$ of random variables indexed by $T>0$, where
\begin{equation}\label{rap_int}
\clC_{T}(u)(\omega):=\,\int_{0}^{T}\,\Big(\clL_{u}(t)(\omega)-\bbE[
\clL_{u}(t)]\Big)\,dt,\qquad \omega\in \tilde \Omega,
\end{equation}
to be the \emph{average $u$-boundary length process}, indexed by $T>0$, associated to $Z$. Note that by stationarity in law of $\clL_u(t)$, we can rewrite $\clC_{T}(u)$ as $\int_{0}^{T}\,\clL_{u}(t)\,dt -\bbE[
\clL_{u}(0)] T$, where $\bbE[
\clL_{u}(0)]$ is as in~\eqref{meanGKF}.

In this paper we are interested in the behaviour of $\clC_T(u)$, as $ T\to +\infty$. To this purpose, in order to take advantage of Wiener-It\^o theory, we first need to check that $\clC_T(u)$ has finite variance.

\begin{lemm}\label{squareL}
For any $T>0$, the random variable $\clC_T(u)$ is square integrable, i.e., $\rVar(\clC_T(u)) < +\infty$.
\end{lemm}

In view of Lemma~\ref{squareL}, whose proof is also postponed to the Appendix~\ref{app-square}, $\clC_T (u)$ in~\eqref{rap_int} can be expanded into so-called Wiener chaoses, by means of the Stroock--Varadhan decomposition, see Section~\ref{sec-Wiener} for details in our setting and~\cite[\S~2.2] {noupebook} for a complete discussion. Briefly, this expansion is based on the fact that the sequence of (normalized) Hermite polynomials
$\{ H_q/\sqrt{q!}\}_{q\,\ge\,0}$
\begin{equation}\label{Herm}
H_0\equiv 1,\qquad H_{q}(u):=(-1)^{q}\phi (u)^{-1}\frac{d^q}{du^q}\phi(u), \ q\ge 1
\end{equation}
(where $\phi$ denotes the probability density function of a standard Gaussian random variable) is a complete orthonormal basis of the space of square integrable functions on the real line with respect to the standard Gaussian measure. (The first polynomials are $H_0(u)=1$, $H_1(u)=u$, $ H_2(u)=u^2-1$, $H_3(u) = u^3 - 3u$.) We can write
\begin{equation}\label{exp_s}
\clC_{T}(u)=\sum_{q=0}^{\infty }\clC_{T}(u)[q],
\end{equation}
where the series is orthogonal and converges in $L^2(\Omega)$, here $\clC_{T}(u)[q]$ denotes the orthogonal projection of $\clC _T(u) $ onto the so-called $q^{\rm th}$ Wiener chaos. In Proposition~\ref{teoexpS} we will determine analytic formulas for these chaotic components. We will exploit the series representation~\eqref{exp_s} to investigate the asymptotic distribution of $\clC_T(u)$ as $T\to +\infty$. Roughly speaking, in the long memory regime the behavior of $\clC_T(u)$ will be determined by a single term of the series, while in the case of short range dependence all chaotic components will contribute in the limit thus influencing both the asymptotic variance and the nature of second order fluctuations of our boundary length functional.

\begin{rema}
This Hermite-type expansion can be directly given for $\clL _u(t)$ in~\eqref{length}, thus being also instrumental for the relations between different geometric processes (evolving over time) associated to the random field $Z$, such as the area of excursion sets and their Euler--Poincaré characteristic, which we plan to investigate in a future paper.
\end{rema}

\subsection{Statement of main results}

Before stating our main results we need some more notation.

\begin{enonce}{Condition}\label{condbeta}
Let $\{ \beta_\ell, \ell \in \widetilde{\bbN}
\}$ be the sequence defined in Assumption~\ref{basic3}.
\begin{itemize}
\item The sequence $\{ \beta_\ell, \ell \in \widetilde{\bbN},
\ell\ge 1\}$ admits minimum. Let us set
\begin{equation}\label{def_star}
\beta_{\ell ^{\star }}:=\min\left\{ \beta _{\ell }, \ell \in \widetilde{
\bbN},\ell\ge 1\right\},\qquad \clI^{\star }:=\left\{\ell\in
\widetilde{\bbN} :\beta _{\ell }=\beta _{\ell ^{\star }}\right\}.
\end{equation}

\item If $\clI^\star\ne \widetilde{\bbN}$, then the sequence $
\{ \beta_\ell, \ell\in \widetilde{\bbN} \setminus \clI ^\star,\ell\ge 1\}$ admits minimum. Let us set
\begin{equation}\label{def_starstar}
\beta_{\ell ^{\star \star }}:=\min \left\{ \beta _{\ell },\ \ell \in \bbN\backslash \clI^{\star },\ell\ge 1\right\}.
\end{equation}
\end{itemize}
\end{enonce}

Note that $\beta_{\ell^\star}, \beta_{\ell^{\star\star}}\in (0,1]$ and for $
\ell\in \clI^\star$, obviously $C_{\ell}(0)> 0$. In words, $\beta _{\ell ^{\star }}$ represents the smallest exponent corresponding to the largest memory, $\clI^{\star }$ the set of multipoles where this minimum is achieved, and $\beta _{\ell ^{\star \star }}$ the second smallest exponent $\beta _{\ell }$ governing the time decay of the autocovariance $ C_\ell$ at some given multipole $\ell$. Note that we are \emph{excluding} the multipole $\ell =0$ by the definition of $\beta_{\ell^\star}$ and $
\beta_{\ell^{\star\star}}$ in~\eqref{def_star} and~\eqref{def_starstar}, on the other hand $\ell=0$ may belong to $\clI^\star$. \emph{From now on we work under Assumption~\ref{condbeta}.}



\subsubsection{Long range dependence}

As briefly anticipated above, for long memory random fields a single chaotic component determines the asymptotic behavior of $\clC_T(u)$. In our setting, the role of dominating term is played by $\clC_T(u)[2]$ (\eqref{exp_s} with $q=2$): we will see in Remark~\ref{2chaos-hermite} that
\begin{multline}\label{caos2}
\clC_T(u)[2]\\
=\frac{\sigma _{1}}{2}\sqrt{\frac{\pi }{2}}\phi (u)\sum_{\ell } \frac{C_{\ell }(0)(2\ell +1)}{4\pi } \left\{ (u^{2}-1)+\frac{
\lambda _{\ell }/2}{\sigma _{1}^{2}}\right\} \int_0^T\int_{\bbS ^{2}}H_{2}\left(\widehat Z_{\ell }(x,t)\right)dxdt,
\end{multline}
where $\lambda_\ell$ still denotes the $\ell^{\rm th}$ eigenvalue of the spherical Laplacian, $\sigma_1$ is defined as in~\eqref{sigma1}, $ H_2(t)=t^2-1$ denotes the second Hermite polynomial, and $\widehat{Z}_\ell$ is defined as
\begin{equation}\label{zhat}
\widehat Z_\ell (x,t) := \frac{Z_\ell(x,t)}{\sqrt{\frac{2\ell+1}{4\pi} C_\ell(0)}}, \quad (x,t)\in \bbS^2\times \bbR
\end{equation}
recalling the content of Section~\ref{sec-rsh}. In particular, the $\widehat Z_\ell$'s are unit variance time-dependent random spherical harmonics.

In this context, it is worth mentioning the work~\cite{LP11} where the authors study quadratic forms of long range dependent Gaussian random fields under general integrability conditions on the spectral density, which are satisfied by a large class of models.

The asymptotic law of $\clC_T(u)[2]$ in~\eqref{caos2} was introduced in~\cite{MRV:20} and it is related to the Rosenblatt distribution, see below.

\begin{defi}
The random variable $X_\beta$ has the standard Rosenblatt distribution (see e.g.~\cite{Ta:75} and also~\cite{DM79,Ta:79}) with parameter $\beta \in (0,
\frac{1}{2})$ if it can be written as
\begin{equation}\label{Xbeta}
X_{\beta}=a(\beta) \int_{(\bbR^{2})^{^{\prime }}}\frac{e^{i(\lambda _{1}+\lambda _{2})}-1}{i(\lambda _{1}+\lambda _{2})}\frac{W(d\lambda _{1})W(d\lambda _{2})}{|\lambda _{1}\lambda _{2}|^{(1-\beta)/2}}\text{,}
\end{equation}
where $W$ is the white noise Gaussian measure on $\bbR$, the stochastic integral is defined in the Ito's sense (excluding the diagonals: as usual, $(\bbR^2)^{\prime }$ stands for the set $\{ (\lambda_1,
\lambda_2)\in \bbR^2:\lambda_1 \ne \lambda_2\}$), and
\begin{align}
a(\beta) &:= \frac{\sigma(\beta)}{2\,\Gamma(\beta)\,\sin\left({(1-\beta)\pi}/{ 2}\right)}\,,\label{abeta}
\\
\intertext{with}
\sigma(\beta) &:= \sqrt{\frac12(1-2\beta)(1-\beta)}\,.\nonumber
\end{align}
Following~\cite{MRV:20}, we say the random vector $V$ satisfies a composite Rosenblatt distribution of degree $N\in \bbN$ with parameters $ c_{1},\dots,c_{N}\in \bbR,$ if
\begin{equation}\label{V}
V=V_{N}(c_{1},\,\dots,\,c_{N};\beta)\mathop{=}^d\sum_{k=1}^{N}c_{k}X_{k\,;\,\beta}
\text{,}
\end{equation}
where $\{ X_{k;\beta}\} _{k=1,\,\dots,\,N}$ is a collection of i.i.d. standard Rosenblatt random variables of parameter $\beta $.
\end{defi}

Note that indeed $\bbE[X_\beta]=0$ and $\Var(X_\beta)=1$. The Rosenblatt distribution was first introduced in~\cite{Ta:75} and has already appeared in the context of spherical isotropic Gaussian random fields as the exact distribution of the correlogramm~\cite{LTT18}. See also~\cite{Leo88, LRT17,VT13}.

Further characterizations of the composite Rosenblatt distribution, for instance in terms of its characteristic function, can be found in~\cite {MRV:20}.

We are now ready to state our first main result. Let us define the standardized average boundary length functional as
\begin{equation}
\widetilde{\clC}_{T}(u):=\frac{\clC_{T}(u)}{\sqrt{\Var\left(\clC_{T}(u)\right)}}.
\end{equation}


\begin{theo}\label{main-theorem}
If $2\beta_{\ell^\star}<\min(\beta_0,1)$, then as $T\to +\infty$,
\begin{equation}\label{caos2dom}
\widetilde{\clC}_{T}(u) =\frac{\clC_{T}(u)[2]}{\sqrt{
\Var\left (\clC_{T}(u)\right) }}+o_{\bbP}(1)\,,
\end{equation}
where $o_{\bbP}(1)$ denotes a sequence converging to zero in probability, $\Var\left(\clC_{T}(u)\right) \sim
\Var\left(\clC_{T}(u)[2]\right)$ and
\begin{multline}\label{asymp_var}
\lim_{T\,\rightarrow\,\infty }\frac{\Var\left(\clC _{T}(u)[2]\right)}{T^{2-2\beta _{\ell^{\star}}}}\\
=\frac{\sigma _{1}^{2}\pi }{4}\phi^{2}(u)\sum_{\ell\,\in\,\clI^\star}\frac{(2\ell+1)^{2}C_
\ell(0)^2}{(1-2\beta_\ell)(1-\beta_\ell)}\left\{\left(u^{2}-1\right)+\frac{\lambda _{\ell}/2}{\sigma_{1}^{2}}\right\}^2\,.
\end{multline}
Assume in addition that $\# \clI^\star$ in~\eqref{def_star} is finite, then as $T\to \infty $, we have that
\begin{equation*}
\widetilde{\clC}_{T}(u)\overset{d}{\longrightarrow}\sum_{\ell\,\in\,\clI^\star} \frac{C_\ell(0)}{\sqrt{v^\star}}\left\{\left(u^{2}-1\right)+\frac{\lambda _{\ell}/2}{\sigma_{1}^{2}}\right\}V_{2\ell+1}\left(1,\,\dots,\,1;\beta_{\ell^\star}\right)\,,
\end{equation*}
where
\begin{equation*}
v^\star=a\left(\beta_{\ell^\star}\right)^2 \sum_{\ell\,\in\,\clI^\star} \frac{2(2\ell+1)^{2} C_\ell(0)^2}{(1-2\beta_\ell)(1-\beta_\ell)} \left\{ \left(u^{2}-1\right)+\frac{\lambda _{\ell}/2}{\sigma_{1}^{2}}\right\}^2\,,
\end{equation*}
$a(\beta_{\ell^\star})$ being as in~\eqref{abeta}.
\end{theo}


To investigate more deeply the structure of the dominating limit variables, we can distinguish between two cases:
\begin{enumerate}
\item \textbf{Non-unique minimum.} This is the situation where at least one multipole has non integrable (over time) autocovariance function, but the cardinality of $\clI^{\star }$ is strictly larger than one, $\#\clI^{\star }>1$, meaning that the minimum of $\{\beta _{\ell }, \ell\in \widetilde{\bbN} \}$ is non-unique. In this case the dominating second order chaos has a neat expression for the variance but the Berry cancellation phenomenon cannot occur, i.e., the variance has the same order of magnitude at any level $u\in\bbR$, see~\eqref{asymp_var}.
\item \textbf{Unique minimum}. This is the situation where at least one multipole has non integrable (over time) autocovariance function and $\#
\clI^{\star }=1,$ meaning that there is a single multipole (labelled $
\ell ^{\star })$ where $\{\beta _{\ell }, \ell\in \widetilde{\bbN }\}$ achieves its minimum, and hence where the temporal dependence is maximal. In these circumstances, we have not only that the boundary length is dominated by the second chaos~\eqref{caos2dom}, but also that this chaos admits an asymptotic expression in terms of the (random) $L^{2}(\bbS ^2) $-norm of the field $Z_{\ell^{\star}}$: from~\eqref{caos2}
\begin{multline}\label{2nd-mono}
\clC_T(u)[2]\\
=\frac{\sigma _{1}}{2}\sqrt{\frac{\pi }{2}}\phi (u)\left(2\ell^{\star }+1\right)\left\{\left(u^{2}-1\right)+\frac{\lambda _{\ell^{\star }}/2}{\sigma _{1}^{2}}\right\} \frac{C_{\ell^\star}(0)}{4\pi} \int_0^T\int_{\bbS ^{2}}H_{2}\left(\widehat Z_{\ell^\star}(x,t)\right)\,dxdt.
\end{multline}
Moreover, from~\eqref{2nd-mono} \emph{perfect correlation occurs between the average boundary length and the average excursion area investigated in~\cite {MRV:20}}. In this case the variance is asymptotic to
\begin{equation}\label{Berry}
\Var\left(\clC_{T}(u)\right)\sim{T^{2-2\beta_{\ell^\star}}}\frac{\sigma_{1}^{2}\pi \phi^{2}(u)}{4} \frac{\left(2\ell^\star+1\right)^{2} C_{\ell^\star}(0)^2}{\left(1-2\beta_{\ell^\star}\right)\left(1-\beta_{\ell^\star}\right)} \left\{\left(u^{2}-1\right)+\frac{\lambda_{\ell^\star}/2}{\sigma_{1}^{2}}\right\}^2 
\end{equation}
\end{enumerate}

\begin{rema}[On the Berry's cancellation phenomenon]
A lot of attention in the last fifteen years has been devoted to the investigation of the so-called \emph{Berry's cancellation phenomenon}; namely, the fact that the variance for the boundary length in the case of random eigenfunctions has been shown to be asymptotically (in the high-energy limit) of lower order in the case of nodal (i.e., zero) sets than for all other thresholds. This fact was first noted in the physical literature by Berry in~\cite{Ber02}; the first rigorous proof for the variance of nodal lines for spherical eigenfunctions is in~\cite{Wig}, whereas the connection with the disappearance of the second order chaos has been given in~\cite{MPRW16, MRW, NPR19} respectively, for eigenfunctions on the torus, on the plane and on the sphere; see also~\cite{CM2018} for extensions to other geometric functionals and~\cite{M2022,Rossi2018, Wig22} for reviews of this area. See also Section~\ref{secHigh}.

\noindent In the case of sphere-cross-time random fields, conditions for the variance of the average excursion area at zero level to be of smaller order were investigated in~\cite{MRV:20}. As far as average boundary lengths is concerned, the picture is actually more subtle; indeed a form of Berry's cancellation phenomenon can hold at non-zero levels, meaning that the variance of boundary lengths at these thresholds has a smaller order of magnitude than the variance at any other level. Indeed, from~\eqref{2nd-mono} and~\eqref{Berry} we see that the leading term in the variance disappears at the points $u$ such that
\begin{equation*}
\left(u^{2}-1\right)+\frac{\lambda _{\ell }/2}{\sigma _{1}^{2}}=0,\qquad\text{ i.e. } u^{2}=1-\frac{\lambda _{\ell ^{\star }}/2}{\sigma _{1}^{2}}\text{.}
\end{equation*}
This clearly implies that Berry's cancellation can occur in a point $\pm u^{\star }\in (0,1)$ such that
\begin{equation*}
u^{\star }=\pm \sqrt{1-\frac{\mathbb{\lambda }_{\ell ^{\star }}}{\bbE _{C}[\mathbb{\lambda }_{\ell }]}},
\end{equation*}
where $\bbE_{C}[\mathbb{\lambda }_{\ell }]$ is the expected value under the probability measure assigning weights $\frac{2\ell +1}{4\pi } C_{\ell }(0)$ (recall that $\sigma_0^2=\sum \frac{2\ell +1}{4\pi }C_{\ell }(0)=1$). In other words, for Berry's cancellation phenomenon to appear we need long memory to occur in multipoles which are ``lower than average'' in terms of the angular power spectrum, or we need the field at any given time to have huge power on low multipoles. Indeed in the standard monochromatic wave case we have $\frac{\mathbb{\lambda }_{\ell ^{\star }}}{\bbE_{C}[\mathbb{
\lambda }_{\ell }]}=1$ and we are back to the nodal length case. At $u^\star$ we have two possible scenarios: if $2\beta_{\ell^{\star\star}}>3\beta_{
\ell^{\star}}$ then the boundary length is dominated by the third chaos, while if $2\beta_{\ell^{\star\star}}<3\beta_{\ell^{\star}}$ then the boundary length is still dominated by its second chaotic component.
\end{rema}

\subsubsection{Short Range Dependence}

If the set of long range dependent multipoles $\clI$ is empty, meaning that $\beta_0=1$ and $2\beta _{\ell }>1$ for all $\ell\ge1$, then the second-order chaotic component would no longer be dominating, and investigation of all terms of the series~\eqref{chaosexpS} below is required. In this case a Gaussian limit via classic Breuer--Major arguments~\cite{BM} holds.

We first need to introduce some more notation: for $q\ge 1$, let
\begin{equation*}
s^2_q:=\lim_{T\,\rightarrow\,\infty}\frac{Var(\clC_T(u)[q])}{T};
\end{equation*}

\begin{theo}\label{quattro}
Assume $\beta_0=1$ and $2\beta _{\ell }>1$ for all $\ell\ge1$. Then we have
\begin{equation*}
\lim_{T\,\rightarrow\,\infty }\frac{\rVar\left(\clC_{T}(u)\right) }{T}=\sum_{q=1}^{+\infty} s_q^2\,,
\end{equation*}
and moreover, as $T\to +\infty$,
\begin{equation*}
\widetilde{\clC}_{T}(u)=\frac{\clC_{T}(u)-\bbE{\clC}_{T}(u)}{\sqrt{Var(\clC_{T}(u))}}\overset{d}{\rightarrow} Z,
\end{equation*}
$Z\sim \clN(0,1)$ being a standard Gaussian random variable.
\end{theo}

Recall that for $\beta_0=1$ we have $\int_{\bbR} C_0(\tau)\,d\tau \in (0,+\infty)$ (see Condition~\ref{basic3}) so that $s_1^2>0$ yielding $
\sum_{q\,\ge\,1}s^2_q >0$ (the limiting variance constant is strictly positive). The proof of the previous result can then be established by a standard (although lengthy) analysis of terms in the chaos expansions~\eqref {chaosexpS}. In particular, note that by the $L^2$ convergence of the Wiener chaoses it is sufficient to focus on an (arbitrarily large but) finite number of components (the remainder may be made negligible, uniformly over $T$); the fourth cumulants of these components can be shown to converge to zero after normalizing for the variance, so that the Central Limit Theorem may follow from Stein--Malliavin arguments (see~\cite{noupebook}). Details are omitted for brevity's sake.



\subsection{Structure of the paper}

In Section~\ref{sec-discussion} we compare our main findings with the existing literature. Section~\ref{secProofsMR} contains the proof of our main theorem, together with the presentation of its main technical tool and it is divided as follows. In Section~\ref{sect4.1} we present the $L^2$ approximation of the length for level curves, whereas the Wiener chaotic decomposition of our boundary length functional is given in Section~\ref{sec-Wiener}; in particular we study the second order chaotic component, i.e. we compute its variance and hence we obtain a much neater asymptotic expression, which includes only the multipoles corresponding to the strongest memory. A much more technical computation is aimed to show that all the higher-order chaotic components are asymptotically negligible; these results are then combined in Section~\ref{sec-Wiener} to prove our main theorem. The Appendix collects a number of important auxiliary results, that derive explicitly covariance structures and cover measurability issues, mean-square approximations and chaotic decompositions.



\section{Discussion}
\label{sec-discussion}

In this Section we compare our main results with the existing related literature on the geometry of random fields.

\subsection{A comparison with the high-energy regime literature}\label{secHigh}

The literature on the geometry of random fields on manifolds has become vast over the last decade, see for instance~\cite{M2022,Rossi2018,Wig22} for some recent surveys. Much of the literature has concentrated on the high-frequency geometry for random eigenfunctions, in the case of random fields on the sphere (or on other Riemannian manifolds, for instance the torus) with no temporal dependence. In particular, concerning level curves it has been shown that the following asymptotic results hold:
\begin{itemize}
\item for level sets corresponding to $u\neq 0,$ the length of level curves is dominated by a single projection term in the chaos expansion, i.e., the second-order component;

\item this component can be expressed in terms of the norm of the function, without its derivatives and it disappears in the nodal case $u=0$ (the so-called Berry's cancellation phenomenon);

\item at $u=0,$ the nodal length is again dominated by a single term in the chaos expansion, which is the fourth-order component;

\item in both cases, it is possible to establish quantitative central limit theorems, in the high-energy limit;

\item the nodal length and the level curves are asymptotically perfectly uncorrelated. However, considering the partial autocorrelation, i.e., removing (or \emph{freezing}) the effect of the random $L^{2}$-norm, the asymptotic correlation is again unity.
\end{itemize}

Much of these results can be extended to other geometric functionals, such as Lipschitz--Killing curvatures (in the two-dimensional case, the excursion area, the boundary length and the Euler--Poincaré characteristic) and critical points. Indeed, full correlation has been shown to hold, in the high-frequency limit, for all these statistics, in \emph{generic} cases where $u\neq 0$ (or, more generally, where the second-order chaos component does not disappear).

The setting we consider in this paper is rather different, for a number of reasons. Firstly, we are not considering eigenfunctions, but arbitrary (although Gaussian and isotropic) spherical random fields. More importantly, we are going beyond those previous results by allowing a form of dependence over time; because of this, the asymptotic framework of our work here is based upon a different asymptotic regime, that is, fluctuations for a growing span over time, rather than for higher and higher frequency eigenfunctions. The results presented here show then both analogies and important differences with the existing literature. More precisely, let us note the following:
\begin{itemize}
\item It is still the case (in the long memory case) that the fluctuations around the expected value are asymptotically (as $T\rightarrow \infty $) dominated by the second order chaos; on one hand, this chaos can again be expressed in terms of the harmonic components of the fields itself, without the need to resort to derivatives (despite the fact that these derivatives do appear in the Kac-Rice representation of level curves, see below).

\item On the other hand, it is no longer the case that the second-order chaos is proportional to the random $L^{2}$-norm of the field itself; it is instead a linear combination of $L^{2}$-norm of the harmonic components $ Z_{\ell }$.

\item Related to the previous point, it is no longer the case that Berry's cancellation occurs at the nodal level $u=0$, and actually for the case $\#
\clI^\star>1$, the Berry cancellation phenomenon cannot occur at all, i.e., the variance has the same order of magnitude at any level $u\in\bbR$.

\item In the case of asymptotically monochromatic fields, i.e., those where the minimum of the memory parameter $\beta_\ell$ is attained on a single multipole, there exist levels where the second-order chaos disappears, and hence the variance is asymptotically of lower order. The exact value of these levels depends upon a combination of the variance of the single component $Z_{\ell^\star}$, and the variance of the derivative of the entire field $Z$. However, rather differently from the literature so far, these do not correspond to the nodal case $u=0$.
\end{itemize}

Let us also recall that in~\cite{MRV:20} the large time behavior of the empirical excursion area of the space-time spherical random field $Z$ in~\eqref{eq1} has been investigated, i.e. the asymptotic distribution of
\begin{equation}\label{area1}
\mathcal{M}_T(u) := \int_0^T \left (\int_{\bbS^2}\left(1_{\left\{ Z(x,t)\ge u\right\}} - \bbP(Z(x,t)\ge u)\right)\,dx\right)\,dt
\end{equation}
as $T\to +\infty$. First of all it is worth mentioning that the analysis of~\eqref{area1} can be carried out under the sole Assumption~\ref{basic}, for the length of level curves instead we need more regularity for $Z$, as explained at the beginning of Section~\ref{sec-MR}. Moreover, for the excursion area the zero-level $u=0$ is still a cancellation point under long memory circumstances, while this is not the case for the variance of level curves. More importantly, for the excursion area the second-order chaos is proportional to the random $L^2$ norm of the random field, whereas for level curves the second-order chaos is proportional to a linear combination of the $L^2$ norms of the eigenfunctions of the field; the two chaoses are hence not perfectly correlated, unless the fields are asymptotically monochromatic.

\section{Proofs of the main results}
\label{secProofsMR}

In this Section we prove our main results. As anticipated in Section~\ref {sec-MR}, the starting point of our argument is the Stroock-Varadhan decomposition of $\clC_T(u)$ in~\eqref{rap_int}.

\subsection{The mean square approximation}\label{sect4.1}

Let $t\in \bbR$, the length of $u$-level curves can be \emph{formally} represented as
\begin{equation*}
\clL_u(t) = \int_{\bbS^2} \delta_u(Z(x,t)) \| \nabla Z(x,t)\|\,dx,
\end{equation*}
where $\delta_u$ is the Dirac mass in $u$. For $\epsilon >0$ consider the $
\epsilon$-approximating $u$-level curves length (see Lemma~\ref{approx2} in Appendix~\ref{app-square})
\begin{equation}\label{Leps}
\clL_u^\epsilon(t) := \frac{1}{2\epsilon}\int_{\bbS^2} 1_{\left[u-\epsilon,u+ \epsilon\right]}(Z(x,t)) \| \nabla Z(x,t)\|\,dx
\end{equation}
and define accordingly the $\epsilon$-approximating random variable
\begin{equation}\label{Ceps}
\clC^\epsilon_T(u) := \int_0^T \big(\clL_u^\epsilon(t) -
\bbE\left[\clL_u^\epsilon(t)\right]\big)\,dt.
\end{equation}
The following technical result is crucial and will be proved in Appendix~\ref {app-square}.

\begin{lemm}\label{approx}
As $\epsilon\to 0$,
\begin{equation}
\clC_T^\epsilon(u) \to \clC_T(u)
\end{equation}
both a.s. and in $L^2(\bbP)$, where $\clC_T(u)$ is defined as in~\eqref{rap_int}.
\end{lemm}

\subsection{The Wiener chaos expansion}
\label{sec-Wiener}

In order to derive the analytic form for~\eqref{exp_s} we get inspired by the chaotic decomposition for level curves of Gaussian random fields found in e.g.~\cite{KL01, MPRW16}. Let us introduce the collection of coefficients $\{\alpha_{n,m} : n,m\geq 1\}$ and $\{\beta_{l}(u) : l\geq 0\}$, related to the (formal) Hermite expansions of the norm $\| \cdot \|$ in $\bbR^2$ and the Dirac mass $\delta_u(\cdot)$ respectively:
\begin{equation}\label{e:beta}
\beta_{q}(u):= \phi(u)H_{q}(u),
\end{equation}
where $\phi$ is the standard Gaussian probability density function, $H_{l}$ denotes the $l^{\rm th}$ Hermite polynomial and $\alpha_{n,m}:=0$ but for the case $n,m$ even
\begin{equation}\label{e:alpha}
\alpha_{2n,2m}:=\sqrt{\frac{\pi}{2}}\frac{(2n)!(2m)!}{n! m!}\frac{1}{2^{n+m}} p_{n+m}\left(\frac14\right),
\end{equation}
where for $N=0, 1, 2,\,\dots$ and $x\in \bbR$
\begin{equation}\label{pN}
p_{N}(x) :=\sum_{j=0}^{N}(-1)^{j}\quad (-1)^{N}{\binom{N}{j}}\ \ \frac{(2j+1)!}{(j!)^2} x^j.
\end{equation}

In view of~\eqref{sigma1}, we define the normalized gradient and derivatives for $j=1,2$
\begin{equation*}
\widetilde \nabla := \nabla /\sigma_1,\quad \widetilde \partial_{j;x} :=
\partial_{j;x}/\sigma_1.
\end{equation*}

\begin{prop}[Chaotic expansion for $\clC_T(u)$]\label{teoexpS}
For every $T>0$ and $q\geq 1$,
\begin{multline} \label{e:ppS}
\clC_T(u)[q] = \sigma_1 \,\sum_{m=0}^{q}\;
\sum_{k=0}^{m} \frac{\alpha _{k,m-k}\beta_{q-m}(u) }{(k)!(m-k)!(q-m)!}\\
\times \int_0^T \int_{\bbS^2} H_{q-m}(Z(x,t)) H_{k}\left(\widetilde\partial_{1;x} Z(x,t)\right)H_{m-k}\left(\widetilde\partial_{2;x} Z(x,t)\right)\,dx\,dt\,.
\end{multline}
As a consequence, one has the representation
\begin{multline}\label{chaosexpS}
\clC_T(u) =\sigma_1\,\sum_{q=1}^{+\infty}
\sum_{m=0}^{q}\sum_{k=0}^{m} \frac{\alpha _{k,m-k}\beta _{q-m}(u)}{ (k)!(m-k)!(q-m)!} \,\\
\times\, \int_0^T\int_{\bbS^2}\!\! H_{q-m}(Z(x,t))H_{k}\left(\widetilde\partial_{1;x} Z(x,t)\right)H_{m-k}\left(\widetilde{\partial} _{2;x} Z(x,t)\right)\,dx\,dt\,,
\end{multline}
where the series converges in $L^2(\Omega)$.
\end{prop}

The proof of Proposition~\ref{teoexpS} is postponed to the Appendix~\ref{app-chaos}: first we compute the chaotic expansion of $\clC^\epsilon_T(u)$ in~\eqref{Ceps}, then we let $\epsilon\to 0$ obtaining~\eqref{chaosexpS} thanks to Lemma~\ref{approx}.

Let us investigate the chaotic components~\eqref{e:ppS} starting from the case $q=1$.

\subsubsection{The first chaotic projection}

The first term in the series~\eqref{chaosexpS} is
\begin{equation}
\begin{split}
\clC_T(u)[1] &= \sigma_1 \sqrt{\frac{\pi}{2}} u\phi(u) \int_{\bbS^2}Y_{0,0}(x)\,dx \int_0^T a_{0,0}(t)\,dt \notag \\
&= \sigma_1 \sqrt 2 \pi u\phi(u) \int_0^T a_{0,0}(t)\,dt.
\end{split}
\end{equation}

\begin{lemm}\label{lemvar1}
We have, as $T\to +\infty$,
\begin{align*}
\lim_{T\,\to\,\infty} \frac{\Var(\clC_T(u)[1])}{T^{2-\beta_0}} &= \sigma_1^2 2\pi^2 u^2 \phi(u)^2\frac{2C_0(0)}{(1-\beta_0)(2-\beta_0)},\quad \text{if } \beta_0\in (0,1)
\\
\intertext{and}
\lim_{T\,\to\,\infty}\frac{\Var(\clC_T(u)[1])}{T} &=
\sigma_1^2 2\pi^2 u^2\phi(u)^2 \int_{\bbR} C_0(\tau)\,d\tau,\quad
\text{if }\beta_0=1.
\end{align*}
\end{lemm}

The proof of Lemma~\ref{lemvar1} is identical to the proof of~\cite[Lem\-ma~4.2]{MRV:20} and hence omitted. Recall that Assumption~\ref{basic3} ensures that $C_0(0)>0$ and that for $\beta_0=1$
\begin{equation*}
\int_{\bbR} C_0(\tau)\,d\tau \in (0,+\infty);
\end{equation*}
it is worth noticing that $\clC_T(u)[1]\ne 0$ if and only if $u\ne 0$.

\subsubsection{The second order chaotic projection}

Recall the notation ($\lambda _{\ell }:=\ell(\ell +1)$) from~\eqref{somma1} and~\eqref{sigma1}
\begin{equation*}
\sigma_{0}^{2} =\sum_{\ell}\frac{(2\ell +1)}{4\pi}C_{\ell}(0)=1,\quad
\sigma_{1}^{2} =\sum_{\ell}\frac{(2\ell +1)}{4\pi}C_{\ell}(0)\frac{\lambda _{\ell}}{2}.
\end{equation*}

\begin{prop}\label{propcaos2}
The second order chaotic component can be written as
\begin{multline}\label{2chaos}
\clC_T(u)[2]\\
=\frac{\sigma _{1}}{2}\sqrt{\frac{\pi }{2}}\phi (u)\sum_{\ell }(2\ell +1)\left\{ \left(u^{2}-1\right)+\frac{\lambda_{\ell }/2}{\sigma _{1}^{2}}
\right\} \int_{[0,T]} \left\{\widehat{C}_{\ell }(t)-C_{\ell }(0)\right\} dt \,,
\end{multline}
where $\widehat{C}_{\ell }(t)$ is the sample power spectrum~\eqref{sampletris}
\begin{equation*}
\widehat{C}_{\ell }(t)=\frac{1}{2\ell +1}\sum_{m=-\ell}^{\ell }\left\vert a_{\ell m}(t)\right\vert^{2}=\frac{1}{2\ell +1}\int_{\bbS^{2}}Z_{\ell}(x,t)^{2}dx\text{.}
\end{equation*}
\end{prop}

Note that for every $t\in \bbR$
\begin{equation*}
\bbE\left[\widehat C_\ell(t)\right] = C_\ell(0).
\end{equation*}


\begin{proof}[{Proof of Proposition~\ref{propcaos2}}]
From Proposition~\ref{teoexpS},
\begin{multline*}
\clC_T(u)[2] = \sigma_1 \,\frac{\alpha _{0,0}\beta _{2}(u)}{2}
\,\int_0^T \int_{\bbS^2} H_{2}(Z(x,t))\,dx\,dt \\
+\sigma_1 \,\frac{\alpha _{2,0}\beta _{0}(u)}{2} \, \int_0^T\int_{\bbS^2}
\left(\left\langle \widetilde \nabla Z(x,t), \widetilde \nabla Z(x,t)\right\rangle - 2 \right)\,dx\,dt.
\end{multline*}
Recall the basic (Green--Stokes) identity for a regular function $T:\bbS ^2\to \bbR$
\begin{equation*}
\int_{\bbS^{2}} \langle \nabla T, \nabla T \rangle dx=-\int_{\bbS ^{2}}T\Delta_{\bbS^2} T\,dx,
\end{equation*}
then
\begingroup
\allowdisplaybreaks
\begin{align*}
&\clC_T(u)[2] = \sigma_1 \frac{\alpha _{00}\beta _{2}(u)}{2}\int_{[0,T]}\int_{\bbS ^{2}} \left(Z(x,t)^{2}-1\right) dxdt
\\
&\quad -\sigma_1 \frac{\alpha _{20}\beta _{0}(u)}{2}\frac{1}{\sigma_1^2}\int_{[0,T]}\int_{
\bbS^{2}} \left(Z(x,t) \Delta_{\bbS^2} Z(x,t) - 2\sigma_1^2\right) d x dt\\
&=\sigma_1 \frac{\alpha _{00}\beta _{2}(u)}{2}\int_{[0,T]}\int_{\bbS ^{2}}\left\{ \left(\sum_{\ell }Z_{\ell }(x,t)\right)^{2}-\sum_{\ell }\frac{ (2\ell +1)}{4\pi }C_{\ell }\right\} dxdt \\
&\quad -\sigma_1 \frac{\alpha _{20}\beta _{0}(u)}{2}\frac{1}{\sigma_1^2}\int_{[0,T]}\int_{\bbS^{2}}\left\{\sum_{\ell }Z_{\ell }(x,t)\Delta \sum_{\ell ^{\prime }}Z_{\ell ^{\prime }}(x,t)-2\sigma_1^2\right\} dx dt\\
&=\sigma_1 \frac{\alpha _{00}\beta _{2}(u)}{2}\int_{[0,T]}\int_{\bbS ^{2}}\left\{\sum_{\ell} \sum_{\ell^{\prime}}Z_{\ell }(x,t)Z_{\ell^{\prime} }(x,t)-\sum_{\ell }\frac{(2\ell +1)}{4\pi }C_{\ell }\right\} dxdt \\
&+\sigma_1 \frac{\alpha _{20}\beta _{0}(u)}{2}\frac{1}{\sigma_1^2}\int_{[0,T]}\int_{\bbS^{2}}\left\{ \sum_{\ell }\sum_{\ell ^{\prime }}Z_{\ell }(x,t)\lambda _{\ell ^{\prime }}Z_{\ell ^{\prime }}(x,t)-\sum_{\ell } \frac{ (2\ell +1)}{4\pi }C_{\ell }\lambda _{\ell }\right\} dxdt\\
&=\sigma_1 \frac{\alpha _{00}\beta _{2}(u)}{2}\int_{[0,T]}\left\{ \sum_{\ell }\int_{\bbS ^{2}} Z_{\ell }(x,t)^{2} dx-4\pi\sum_{\ell }\frac{(2\ell +1)}{4\pi }C_{\ell }\right\} dt \\
&+\sigma_1 \frac{\alpha _{20}\beta _{0}(u)}{2}\frac{1}{\sigma_1^2}
\int_{[0,T]}\left\{\sum_{\ell }\lambda _{\ell }\int_{\bbS^{2}}Z_{\ell }(x,t)^2 dx-4\pi \sum_{\ell } \lambda _{\ell } \frac{(2\ell +1)}{4\pi }C_{\ell } \right\} dt\\
&=\sigma_1 \frac{\alpha _{00}\beta _{2}(u)}{2}\int_{[0,T]}\left\{ \sum_{\ell}(2\ell+1)\widehat{C}_{\ell }(t)-\sum_{\ell}(2\ell +1)C_{\ell}\right\} dt
\\
&+\sigma_1 \frac{\alpha _{20}\beta _{0}(u)}{2}\frac{1}{\sigma_1^2}
\int_{[0,T]}\left\{ \sum_{\ell }\lambda _{\ell }(2\ell+1)\widehat{C}_{\ell }(t)- \sum_{\ell} \lambda _{\ell} (2\ell +1)C_{\ell} \right\} dt \\
&=\sigma_1 \frac{\alpha _{00}\beta _{2}(u)}{2}\int_{[0,T]}\sum_{\ell}(2\ell +1)\left\{ \widehat{C}_{\ell }(t)-C_\ell\right\} dt \\
&+\sigma_1 \frac{\alpha _{20}\beta _{0}(u)}{2}\frac{1}{\sigma_1^2}
\int_{[0,T]}\sum_{\ell }(2\ell +1)\ell (\ell +1)\left\{\widehat{C}_{\ell }(t)-C_\ell\right\}dt,
\end{align*}
\endgroup
where $\widehat{C}_{\ell }$ is given by
\begin{equation}\label{sampletris}
\widehat{C}_{\ell }(t):=\frac{1}{2\ell +1}\sum_{m=-\ell }^{\ell }\left\vert a_{\ell m}(t)\right\vert^{2}=\frac{1}{2\ell +1}\int_{\bbS^{2}}Z_{\ell}(x,t)^{2}dx\text{.}
\end{equation}
Also,
\begin{equation*}
\beta _{0}(u)=\phi (u)\text{, }\beta _{2}(u)=\phi (u)\left(u^{2}-1\right)\text{, }
\alpha _{00}=\sqrt{\frac{\pi }{2}}\text{, }\alpha _{02}=\frac{1}{2}\sqrt{
\frac{\pi }{2}}
\end{equation*}
whence
\begin{equation*}
\frac{\alpha _{00}\beta _{2}(u)}{2}=\frac{1}{2}\sqrt{\frac{\pi }{2}}\phi (u)\left(u^{2}-1\right)\text{, }\frac{\alpha _{20}\beta _{0}(u)}{2}=\frac{1}{4}\sqrt{
\frac{\pi }{2}}\phi (u)\text{.}
\end{equation*}
We can then write the second-order chaos more compactly as
\begin{align*}
\clC_T(u)[2] &= \frac{\sigma _{1}}{2}\sqrt{\frac{\pi }{2}}\phi (u)\left(u^{2}-1\right)\int_{[0,T]}\sum_{\ell }(2\ell +1)\left\{ \widehat{C}_{\ell }(t)-C_{\ell }\right\} dt \\
&\quad+\frac{\sigma _{1}}{4}\sqrt{\frac{\pi }{2}}\phi (u)\frac{1}{\sigma _{1}^{2} }\int_{[0,T]}\sum_{\ell }(2\ell +1)\lambda _{\ell }\left\{ \widehat{C}_{\ell }(t)-C_{\ell }\right\} dt \\
&= \frac{\sigma _{1}}{2}\sqrt{\frac{\pi }{2}}\phi (u)\sum_{\ell }(2\ell +1)\left\{\left(u^{2}-1\right)+\frac{\lambda _{\ell }/2}{\sigma _{1}^{2}}\right\}
\int_{[0,T]} \left\{\widehat{C}_{\ell }(t)-C_{\ell }\right\} dt
\end{align*}
thus concluding the proof of Proposition~\ref{propcaos2}.
\end{proof}

\begin{rema}\label{2chaos-hermite}
The second order chaos can be also written in terms of Hermite polynomials, since
\begin{multline*}
\int_{[0,T]}\left\{ \widehat{C}_{\ell }(t)-\bbE\widehat{C}_{\ell }(t)\right\} dt \\
\begin{aligned}
&= \int_{[0,T]}\left\{\frac{1}{2\ell+1}\int_{\bbS^2}Z_\ell(x,t)^2dx-
\frac{1}{2\ell+1}\int_{\bbS^2}\bbE\left[Z_\ell(x,t)^2\right]\right\} dt \\
&= \int_{[0,T]}\left\{\frac{1}{2\ell+1}\int_{\bbS^2}Z_\ell(x,t)^2dx-
\frac{1}{2\ell+1}\int_{\bbS^2}\frac{2\ell+1}{4\pi}C_{\ell }(0)dx\right\} dt \\
&= \int_{[0,T]}\frac{1}{2\ell+1}\int_{\bbS^2}\left\{Z_\ell(x,t)^2dx-
\frac{2\ell+1}{4\pi}C_{\ell }(0)\right\} dxdt \\
&= \int_{[0,T]}\frac{1}{2\ell+1}\frac{2\ell+1}{4\pi}C_{\ell }(0)\int_{
\bbS^2}\left\{\widehat Z_\ell(x,t)^2dx-1\right\}dxdt\,,
\end{aligned}
\end{multline*}
so that
\begin{multline*}
\clC_T(u)[2]=\\
\frac{\sigma _{1}}{2}\sqrt{\frac{\pi }{2}}\phi (u)\sum_{\ell } \frac{C_{\ell }(0)(2\ell +1)}{4\pi } \left\{\left(u^{2}-1\right)+\frac{\lambda _{\ell }/2}{\sigma _{1}^{2}}\right\} \int_0^T\int_{\bbS^{2}}H_{2}\left(\widehat Z_{\ell }(x,t)\right)dxdt\text{,}
\end{multline*}
as anticipated in~\eqref{caos2}.
\end{rema}

\begin{rema}[Non-asymptotic monochromatic field]\label{2chaos_mono}
In the special case of monochromatic fields where
\begin{equation*}
C_{\ell }(0)\neq 0\Leftrightarrow \ell =\ell ^{\star },
\end{equation*}
we have that $(\sigma^2_0= \frac{2\ell+1}{4\pi}C_\ell(0)=1)
\sigma^2 _{1}=\frac{\ell (\ell +1)}{2}$ and we get a straightforward generalization of the standard non-asymptotic expression for the second-order chaos for the boundary length of a time-dependent random spherical harmonic, namely
\[
\clC_T(u)[2] =\sqrt{\frac{\ell (\ell +1)}{2}}\frac{1}{2}\sqrt{\frac{
\pi }{2}}u^{2}\phi (u)\int_{[0,T]}\int_{\bbS^{2}} H_2(Z_\ell(x,t)) dxdt.
\]
(Note that $\widehat{Z}_\ell = Z_\ell$ in this case.)
\end{rema}

\begin{rema}
It is clear from~\eqref{2chaos} that the disappearance of the second-order chaos at $u=0$ (closely related to the Berry's cancellation phenomenon) does not occur for non-monochromatic space-time random fields -- although it does occur in the non-asymptotic monochromatic case (Remark~\ref {2chaos_mono}). As we already showed in Section~\ref{sec-MR}, the cancellation can occur asymptotically (as $T\rightarrow \infty $) in some cases of long range dependent fields where the memory parameter attains its minimum on a single multipole $\ell^{\star}$; this can be viewed as a form of asymptotic monochromatic behaviour.
\end{rema}

From~\eqref{caos2} we have
\begingroup
\allowdisplaybreaks
\begin{align}\label{varlunga2}
&\Var\left(\clC_{T}(u)[2]\right)\\
&=\frac{\sigma _{1}^{2}\pi }{8}\phi ^{2}(u)\bbE\left[\left(
\sum_{\ell }(2\ell +1)\left\{\left(u^{2}-1\right)+\frac{\lambda _{\ell }/2}{\sigma _{1}^{2}}\right\}\right.\right.\nonumber\\
&\mkern200mu\left.\left.\frac{C_{\ell }(0)}{4\pi }\int_{[0,T]}\int_{\bbS^{2}}H_{2}\left(\widehat{Z}_{\ell }(x,t)\right)dxdt\right) ^{2}\right]\nonumber\\
&=\frac{\sigma _{1}^{2}\pi }{8}\phi ^{2}(u)\sum_{\ell }(2\ell +1)^{2}\left\{ (u^{2}-1)+\frac{\lambda _{\ell }/2}{\sigma _{1}^{2}}\right\} ^{2}\frac{C_{\ell }(0)^{2}}{(4\pi)^{2}}\bbE\nonumber\\
&\mkern230mu\left[\left(
\int_{[0,T]}\int_{\bbS^{2}}H_{2}\left(\widehat{Z}_{\ell }(x,t)\right)dxdt\right) ^{2}\right]\nonumber\\
&=\frac{\sigma _{1}^{2}\pi }{8}\phi ^{2}(u)\sum_{\ell }(2\ell +1)^{2}\left\{\left(u^{2}-1\right)+\frac{\lambda _{\ell }/2}{\sigma _{1}^{2}}\right\} ^{2}\frac{C_{\ell }(0)^{2}}{(4\pi)^{2}}\nonumber\\
&\qquad\times \int_{[0,T]^{2}}\int_{\bbS^{2}\times \bbS^{2}}\bbE\left[H_{2}\left(\widehat{Z}_{\ell }(x,t)\right)H_{2}\left(\widehat{Z}_{\ell }(y,s)\right)\right] dxdtdyds\nonumber\\
&=\frac{\sigma _{1}^{2}\pi }{8}\phi ^{2}(u)\sum_{\ell }\left\{\left(u^{2}-1\right)+\frac{\lambda _{\ell }/2}{\sigma _{1}^{2}}\right\} ^{2}\nonumber\\
&\mkern150mu\int_{[0,T]^{2}}\int_{\bbS^{2}\times \bbS^{2}}2\,\frac{(2\ell +1)^{2}}{(4\pi)^{2}}C_{\ell }(t-s)^{2}P_{\ell }(\langle x,y\rangle)^{2}dxdtdyds\nonumber\\
&=\frac{\sigma _{1}^{2}\pi }{8}\phi ^{2}(u)\sum_{\ell }\frac{2\,(2\ell +1)^{2}}{(4\pi)^{2}}\left\{\left(u^{2}-1\right)+\frac{\lambda _{\ell }/2}{\sigma _{1}^{2}}\right\} ^{2}\nonumber\\
&\mkern150mu\int_{[0,T]^{2}}C_{\ell }(t-s)^{2}dtds\int_{\bbS^{2}\times \bbS^{2}}P_{\ell }(\langle x,y\rangle)^{2}dxdy\nonumber\\
&=\frac{\sigma _{1}^{2}\pi }{4}\phi ^{2}(u)\sum_{\ell =0}^{+\infty}(2\ell +1)^{2}\left\{ (u^{2}-1)+\frac{\lambda _{\ell }/2}{\sigma _{1}^{2}}\right\} ^{2}\int_{[0,T]^{2}}C_{\ell }(t-s)^{2}\,dt\,ds,\nonumber
\end{align}
\endgroup
\noindent recalling that $\int_{\bbS^{2}\times \bbS^{2}}P_{\ell }(\langle x,y\rangle)^{2}dxdy = (4\pi)^2/(2\ell+1)$. In view of~\eqref{varlunga2}, we will need the following result.

\begin{lemm}[{\cite[Lem\-ma~4.3]{MRV:20}}]\label{lemma-var2nd}
Fix $\ell \in \widetilde{ \bbN}$. If $2\beta _{\ell}<1$, then
\begin{equation*}
\lim_{T\,\rightarrow\,\infty}\frac{1}{T^{2-2\beta _{\ell }}}
\int_{[0,T]^2}C_{\ell }^{2}(t-s)dtds=\frac{2C_{\ell }(0)^{2}}{ (1-\beta_{\ell})(1-2\beta_{\ell})}\text{.}
\end{equation*}
If $2\beta _{\ell}>1$, then
\begin{equation*}
\lim_{T\,\rightarrow\,\infty}\frac{1}{T}\int_{[0,T]^2}C_{\ell }^{2}(t-s)dtds=\int_{\bbR}C_\ell(\tau)^2\, d\tau\,.
\end{equation*}
\end{lemm}

\begin{prop}\label{prop:var-2nd-chaos}
For $2\beta_{\ell^*}<1$ and $\beta_0\le
\beta_{\ell^*}$ we have that
\begin{multline}
\lim_{T\,\rightarrow\,\infty }\frac{\Var\left(\clC _{T}(u)[2]\right) }{T^{2-2\beta _{\ell^{\star}}}}\\
=\frac{\sigma _{1}^{2}\pi }{4}\phi^{2}(u)\sum_{\ell\,\in\,\clI^\star}\frac{(2\ell+1)^{2}C_
\ell(0)^2}{(1-2\beta_\ell)(1-\beta_\ell)}\left\{\left(u^{2}-1\right)+\frac{\lambda _{\ell}/2}{\sigma_{1}^{2}}\right\}^2\,.
\end{multline}
For $2\beta_{\ell^*} >1$ and $2\beta_0 >1$ we have that
\begin{multline}\label{T2}
\lim_{T\,\rightarrow\,\infty }\frac{\Var\left(\clC _{T}(u)[2]\right) }{T}
\\= \frac{\sigma _{1}^{2}\pi}{4}\phi ^{2}(u)\sum_{\ell=0}^{\infty}(2\ell+1)^{2} \left\{ (u^{2}-1)+\frac{\lambda _{\ell}/2}{\sigma_{1}^{2}}\right\}^2 \int_{(-\infty, +\infty)}C_{\ell }^{2}(\tau)d\tau.
\end{multline}
\end{prop}

\begin{proof}
For $2\beta_{\ell^*}<1$ and $\beta_0\le \beta_{\ell^*}$, from Lemma~\ref {lemma-var2nd}, bearing in mind~\eqref{eqConv}, we can use Dominated Convergence Theorem (as well as~\cite[Lemmas~4.7 and~4.8]{MRV:20}) to get
\begin{align*}
&\lim_{T\,\to\,\infty}\frac{\Var\left(\clC_T (u)[2]\right)}{T^{2-2\beta_{\ell^\star}}}\\
&= \lim_{T\,\to\,\infty} \frac{\sigma _{1}^{2}\pi }{4}\phi ^{2}(u)\sum_{\ell=0}^{\infty}\frac{T^{2-2\beta_{\ell}}(2\ell+1)^{2}}{T^{2-2\beta _{\ell^\star}}}\left\{\left(u^{2}-1\right)+\frac{\lambda _{\ell}/2}{\sigma_{1}^{2}}\right\}^2 \int_{[0,T]^2}\frac{C_{\ell }^{2}(t-s)}{T^{2-2\beta _{\ell }}}dtds\\
&= \frac{\sigma _{1}^{2}\pi }{4}\phi ^{2}(u)\sum_{\ell=0}^{\infty}\lim_{T\,\to\,\infty} \frac{T^{2-2\beta_{\ell}}(2\ell+1)^{2}}{T^{2-2\beta _{\ell^\star}}}\left\{\left(u^{2}-1\right)+\frac{\lambda _{\ell}/2}{\sigma_{1}^{2}}\right\}^2 \frac{C_{\ell }(0)^{2}}{(1-\beta_{\ell})(1-2\beta_{\ell})}dtds\\
&= \frac{\sigma _{1}^{2}\pi }{4}\phi ^{2}(u)\sum_{\ell\,\in\,\clI^\star}\frac{(2\ell+1)^{2}C_\ell(0)^2}{(1-2\beta_\ell)(1-\beta_\ell)}\left\{ \left(u^{2}-1\right)+\frac{\lambda _{\ell}/2}{\sigma_{1}^{2}}\right\}^2.
\end{align*}
The proof of~\eqref{T2} is analogous and hence omitted.
\end{proof}

\subsubsection{Higher-order chaotic projections}
\label{high}

Let us investigate the asymptotic distribution, as $T\to +\infty$, of $
\clC_T(u)$ for $q\ge 3$. In the short memory case, it is trivial to see that, as $T\to +\infty$,
\begin{equation*}
\text{Var}(\clC_T(u)[q]) = O(T).
\end{equation*}
Note that the constants involved in the bound depend on $q$, but they are uniformly square summable. The terms which are of smaller order are clearly negligible; it is thus sufficient to establish that the fourth order cumulants of the non-negligible chaotic components are $o(T^2)$. The proof of this upper bound is standard and straightforward, following the same steps as given for instance in~\cite{MRV:20}.

The next Proposition refers to the long memory case and shows that all chaotic components other than the leading one are uniformly negligible, in the limit $T\to +\infty$.

\begin{prop}\label{var-higher-chaoses}
For $2\beta_{\ell^*} < \min\{ \beta_0, 1\}$, as $T\to +\infty$,
\begin{equation*}
\sum_{q\,\geq\,3}\Var(\clC_{T}(u)[q])=O\left(T^{2-\frac52
\beta _{\ell ^{\star }}}\right).
\end{equation*}
\end{prop}

\begin{proof}
We have
\begin{multline*}
\sum_{q\,\geq\,3}\Var(\clC_{T}(u)[q]) \\
\begin{aligned}
=&\sigma _{1}^{2}\sum_{q\,\geq\,3}\Var\left(
\sum_{m=0}^{q}\sum_{k=0}^{m}\frac{\alpha _{k,m-k}\beta _{q-m}(u)}{k!(m-k)!(q-m)!}\right.\\
&\qquad\left.\int_{0}^{T}\int_{\bbS^{2}}H_{q-m}(Z(x,t))H_{k}\left(
\widetilde{\partial }_{1,x}Z(x,t)\right)H_{m-k}\left(\widetilde{\partial } _{2,x}Z(x,t)\right)dxdt\right) \\
=&\sigma _{1}^{2}\sum_{q\,\geq\,3}\;\sum_{m_{1}=0}^{q}\;\sum_{k_{1}=0}^{m_{1}}\;\sum_{m_{2}=0}^{q}\;
\sum_{k_{2}=0}^{m_{2}}\frac{\alpha _{k_{1},m_{1}-k_{1}}\beta _{q-m_{1}}(u)}{ k_{1}!(m_{1}-k_{1})!(q-m_{1})!}\frac{\alpha _{k_{2},m_{2}-k_{2}}\beta _{q-m_{2}}(u)}{k_{2}!\left(m_{2}-k_{2}\right)!\left(q-m_{2}\right)!} \\
& \int_{[0,T]^{2}}\int_{\bbS^{2}\times \bbS^{2}}
\bbE\bigg[H_{q-m_{1}}(Z(x,t))H_{k_{1}}\left(\widetilde{\partial } _{1,x}Z(x,t)\right)H_{m_{1}-k_{1}}\left(\widetilde{\partial }_{2,x}Z(x,t)\right) \\
&\qquad H_{q-m_{2}}(Z(y,s))H_{k_{2}}\left(\widetilde{\partial}_{1,x}Z(y,s)\right)H_{m_{2}-k_{2}}\left(\widetilde{\partial}_{2,x}Z(y,s)\right)
\bigg] dxdydtds.
\end{aligned}
\end{multline*}
Hence we can write
\begin{multline*}
\sum_{q\,\geq\,3}\Var\big(\clC_{T}(u)[q]\big)\\
\leq \sigma _{1}^{2}\sum_{q\,\geq\,3}\;\sum_{i_{1}+i_{2}+i_{3}=q}\;\sum_{j_{1}+j_{2}+j_{3}=q}\frac{\left|\alpha _{i_{1},i_{2}}\beta _{i_{3}}(u)\right|}{i_{1}!i_{2}!i_{3}!}\frac{\left|\alpha _{j_{1},j_{2}}\beta _{j_{3}}(u)\right|}{j_{1}!j_{2}!j_{3}!}
\,U_{q}\left(i_{1},i_{2},i_{3},j_{1},j_{2},j_{3}\right)\,,
\end{multline*}
where $U_{q}(i_{1},i_{2},i_{3},j_{1},j_{2},j_{3})$ is a sum of at most $q!$ terms of the type
\begin{equation}
\int_{[0,T]^{2}}\int_{\bbS^{2}\times \bbS ^{2}}\prod_{u=1}^{q}\bbE\left[\widetilde{\partial }_{l_{u},x}Z(x,t)
\widetilde{\partial }_{h_{u},x}Z(y,s)\right] \,dx\,dy\,dt\,ds\,,
\end{equation}
where $l_{u},h_{u}\in \{0,1,2\}$ and by $\widetilde{\partial } _{l_{u},x}Z(x,t)$ we denote the normalized partial derivatives with respect to the first or second variable (in our convention, if $l_u=0$ then $
\widetilde{\partial }_{0,x}Z(x,t)=Z(x,t)$). In particular,

\begin{multline*}
\vspace*{-10pt}
\left|U_{q}\left(i_{1},i_{2},i_{3},j_{1},j_{2},j_{3}\right)\right|\\
\leq q!\,(4\pi)^{2}\,\sum_{\ell _{1},\,\dots,\,\ell _{q}=0}^{\infty }\frac{(2\ell _{1}+1)\ell _{1}^{2}}{4\pi}\frac{(2\ell _{2}+1)\ell _{2}^{2}}{4\pi}\cdots \frac{(2\ell _{q}+1)\ell _{q}^{2}}{4\pi} \\
\frac{C_{\ell_1}(0)C_{\ell_2}(0)\cdots C_{\ell_q}(0)\,T^{2-\left(\beta_{\ell_1}+\beta_{\ell_2}+\cdots+\beta_{
\ell_q}\right)}}{\left(1-\beta_{\ell_1}-\beta_{\ell_2}-\cdots-\beta_{\ell_q}\right)\left(2-\beta_{\ell_1}-\beta_{\ell_2}-\cdots-\beta_{\ell_q}\right)}+O(T)\\
\leq \frac{q!\,(4\pi)^{2}\,T^{2-q\beta_{\ell ^{\star }}}}
{\min_{q}\left\{\left(1-\beta_{\ell_{1}}-\beta_{\ell _{2}}-\cdots -\beta_{\ell_{q}}\right)\left(2-\beta_{\ell _{1}}-\beta_{\ell_{2}}-\cdots -\beta_{\ell_{q}}\right)\right\}}
\\
\left(\sum_{\ell =0}^{\infty }\frac{C_\ell(0)(2\ell +1)\ell^{2}}{4\pi}\right)^{q}\,.
\end{multline*}
Indeed, let us investigate the behavior of one of these terms:
\begingroup
\allowdisplaybreaks
\begin{align*}
\int_{[0,T]^2}&\int_{\bbS^2 \times \bbS^2}
\bbE\left[\partial_{2,x} Z(x,t)\partial_{2,y} Z(y,s) \right]^2 \bbE\left[Z(x,t)\partial_{2,y} Z(y,s)
\right]^{q-2}\,dx\,dy\,dt\,ds\\
& =\int_{[0,T]^2}\;\int_{\bbS^2 \times \bbS ^2} \left(\sum_{\ell_1=0}^\infty C_{\ell_1}(t-s)\frac{(2\ell_1+1)}{4\pi}
\partial_{2;x}\partial_{2\,;\,y}P_{\ell_1}(\langle x,y \rangle) \right)^2
\,\\
&\qquad\qquad \times\, \left(\sum_{\ell_3=0}^\infty C_{\ell_3}(t-s)\frac{ (2\ell_3+1)}{4\pi}\partial_{2\,;\,y}P_{\ell_3}(\langle x,y \rangle) \right)^{q-2}\\ 
& = \int_{[0,T]^2}\;\int_{\bbS^2 \times \bbS ^2} \sum_{\ell_1,\ell_2=0}^\infty C_{\ell_1}(t-s)C_{\ell_2}(t-s)\frac{(2\ell_1+1)}{4\pi}\frac{(2\ell_2+1)}{4\pi}\\
&\qquad\qquad\qquad \partial_{2;x}\partial_{2;y}P_{\ell_1}(\langle x,y \rangle)\partial_{2\,;\,x}\partial_{2;y}P_{\ell_2}(\langle x,y \rangle) \\
&\qquad\qquad \sum_{\ell_3,\dots,\ell_q=0}^\infty C_{\ell_3}(t-s)\cdots C_{\ell_q}(t-s)\frac{(2\ell_3+1)}{4\pi}\cdots\frac{(2\ell_q+1)}{4\pi}\\
&\qquad\qquad\qquad
\partial_{2\,;\,y}P_{\ell_3}(\langle x,y \rangle) \cdots
\partial_{2\,;\,y}P_{\ell_q}(\langle x,y \rangle) \\
& \leq(4\pi)^2\sigma _{1}^{2} \int_{[0,T]^2}
\sum_{\ell_1,\ell_2=0}^\infty C_{\ell_1}(t-s)C_{\ell_2}(t-s)\frac{(2\ell_1+1) }{4\pi}\frac{(2\ell_2+1)}{4\pi}\,\ell_1^2 \ell_2^2 \\
&\qquad\qquad \sum_{\ell_3,\,\dots,\,\ell_q=0}^\infty C_{\ell_3}(t-s)\cdots C_{\ell_q}(t-s)\frac{(2\ell_3+1)}{4\pi}\cdots\frac{(2\ell_q+1)}{4\pi}\,
\ell_3 \cdots \ell_q \\
& = \sum_{\ell_1,\,\dots,\,\ell_q=0}^\infty\frac{ (2\ell_1+1)}{4\pi}\frac{(2\ell_2+1)}{4\pi}\,\ell_1^2 \ell_2^2 \frac{ (2\ell_3+1)}{4\pi}\cdots\frac{(2\ell_q+1)}{4\pi}\, \ell_3 \cdots \ell_q
\\
&\qquad \int_{[0,T]^2} C_{\ell_1}(t-s)\cdots C_{\ell_q}(t-s)\,dt\,ds \\
&= \sum_{\ell_1,\,\dots,\,\ell_q=0}^\infty\frac{ (2\ell_1+1)}{4\pi}\frac{(2\ell_2+1)}{4\pi}\, \frac{(2\ell_3+1)}{4\pi}\cdots
\frac{(2\ell_q+1)}{4\pi} \ell_1^2 \ell_2^2 \, \ell_3 \cdots \ell_q
\\
&\qquad\qquad\qquad\frac{C_{\ell_1}(0)C_{\ell_2}(0)\cdots C_{\ell_q}(0)\,T^{2-\left(\beta_{\ell_1}+\beta_{\ell_2}+\cdots+\beta_{
\ell_q}\right)}}{\left(1-\beta_{\ell_1}-\beta_{\ell_2}-\cdots-\beta_{\ell_q}\right)\left(2-\beta_{
\ell_1}-\beta_{\ell_2}-\cdots-\beta_{\ell_q}\right)}+O(T)\,,
\end{align*}
\endgroup
where for the last equality we used~\cite[Lemma~4.11]{MRV:20} and the fact that, from~\eqref{eqConv}, $\sum_{\ell=0}^{\infty }\frac{C_{\ell }(0)(2\ell +1)\ell^2}{4\pi } <+\infty $. As a consequence,
\begingroup
\allowdisplaybreaks
\begin{align*}
&\sum_{q\,\geq\,3}\Var(\clC_{T}(u)[q])
\\
&\leq \sigma_{1}^{2}\sum_{q\,\geq\,3}\sum_{i_{1}+i_{2}+i_{3}=q}\sum_{j_{1}+j_{2}+j_{3}=q}\frac{\left|\alpha_{i_{1},i_{2}}\beta _{i_{3}}(u)\right|}{i_{1}!i_{2}!i_{3}!}\frac{\left|\alpha _{j_{1},j_{2}}\beta_{j_{3}}(u)\right|}{j_{1}!j_{2}!j_{3}!}
\\
&\,\frac{q!\,(4\pi)^{2}\,T^{2-q\beta_{\ell^{\star }}}}{
\min_{q}\left\{\left(1-\beta_{\ell_{1}}-\beta_{\ell_{2}}-\cdots -\beta_{\ell_{q}}\right)\left(2-\beta_{\ell _{1}}-\beta_{\ell_{2}}-\cdots -\beta_{\ell_{q}}\right)\right\}}
\\
& \mkern250mu\left(\sum_{\ell =0}^{\infty }\frac{C_\ell(0)(2\ell +1)\ell^{2}}{4\pi}\right)^{q}
\\
&=\sigma_{1}^{2}\sum_{q\,\geq\,3}\frac{q!\,(4\pi)^{2}\,T^{2-q\beta_{\ell^{\star }}}}{\min_{q}\left\{\left(1-\beta_{\ell _{1}}-\beta_{\ell _{2}}-\cdots -\beta_{\ell_{q}}\right)\left(2-\beta_{\ell_{1}}-\beta_{\ell _{2}}-\cdots -\beta_{\ell_{q}}\right)\right\}}
\\
&\qquad \left(\sum_{i_{1}+i_{2}+i_{3}=q}\frac{\left|\alpha _{i_{1},i_{2}}\beta _{i_{3}}(u)\right|}{i_{1}!i_{2}!i_{3}!}\left(\sum_{\ell =0}^{\infty }\frac{C_\ell(0)(2\ell +1)\ell ^{2}}{4\pi}\right)^{\frac{i_{1}+i_{2}+i_{3}}{2}}\right) ^{2} \\
&=\frac{\sigma_{1}^{2}\,(4\pi)^{2}\,T^{2-\frac52\beta _{\ell^{\star }}}}{\min_{q}\left\{\left(1-\beta_{\ell _{1}}-\beta_{\ell_{2}}-\cdots -\beta_{\ell_{q}}\right)\left(2-\beta_{\ell_{1}}-\beta _{\ell_{2}}-\cdots -\beta_{\ell_{q}}\right)\right\}}
\\
&\qquad \sum_{q\,\geq\,3}q!\,\left(\sum_{i_{1}+i_{2}+i_{3}=q}\frac{\left|\alpha_{i_{1},i_{2}}\beta_{i_{3}}(u)\right|}{i_{1}!i_{2}!i_{3}!}\left(\sum_{\ell =0}^{\infty}\frac{C_\ell(0)(2\ell +1)\ell^{2}}{4\pi \, T^{\left(1-\frac5{2q}\right)\beta_{\ell^{\star}}}}\right)^{\frac{ i_{1}+i_{2}+i_{3}}{2}}\right)^{2}\,.
\end{align*}
\endgroup
So that for each $\varepsilon>0$ there exists $T_\varepsilon>0$ such that
\[
\left(\sum_{\ell =0}^{\infty }\frac{C_\ell(0)(2\ell +1)\ell ^{2}}{4\pi \, T^{\left(1-\frac5{2q}\right)\beta_{\ell^{\star}}}}\right)^{\frac{ i_{1}+i_{2}+i_{3}}{2}} =\left(\sum_{\ell =0}^{\infty }\frac{ C_\ell(0)\,(2\ell +1)\ell ^{2}}{4\pi \, T^{\left(1-\frac5{2q}\right)\beta_{\ell^{\star}}}}\right)^{q/2}<
\varepsilon^{q/2} \,,
\]
for each $T \ge T_\varepsilon$. Hence
\begingroup
\allowdisplaybreaks
\begin{align*}
\sum_{q\,\geq\,3}q! &\left(\sum_{i_{1}+i_{2}+i_{3}=q}\frac{\left|\alpha _{i_{1},i_{2}}\beta _{i_{3}}(u)\right|}{i_{1}!i_{2}!i_{3}!}\left(\sum_{\ell =0}^{\infty}\frac{C_\ell(0)\,(2\ell +1)\ell ^{2}}{4\pi \, T^{\left(1-\frac5{2q}\right)\beta_{\ell^{\star}}}}\right)^{\frac{ i_{1}+i_{2}+i_{3}}{2}}\right)^{2}
\\
&=\sum_{q\,\geq\,3}q!\sum_{i_{1}+i_{2}+i_{3}=q}\sum_{j_{1}+j_{2}+j_{3}=q}\frac{\left|\alpha _{i_{1},i_{2}}\beta _{i_{3}}(u)\right|}{i_{1}!i_{2}!i_{3}!}\frac{\left|\alpha _{j_{1},j_{2}}\beta _{j_{3}}(u)\right|}{j_{1}!j_{2}!j_{3}!}\\
&\quad\times\left(\sum_{\ell =0}^{\infty }\frac{C_\ell(0)\,(2\ell +1)\ell^{2}}{4\pi \, T^{\left(1-\frac5{2q}\right)\beta_{\ell^{\star}}}}\right)^{q}
\\
&\leq \sum_{q\,\geq\,3}q!\sum_{i_{1}+i_{2}+i_{3}=q}\sum_{j_{1}+j_{2}+j_{3}=q}
\frac{\left|\alpha _{i_{1},i_{2}}\beta _{i_{3}}(u)\right|}{i_{1}!i_{2}!i_{3}!}\frac{\left|\alpha _{j_{1},j_{2}}\beta _{j_{3}}(u)\right|}{j_{1}!j_{2}!j_{3}!}\varepsilon^{q}.
\end{align*}
\endgroup
Now, arguing as in~\cite[Section~6.2.2]{DNPR:16}, we have that the previous quantity is equal to
\begingroup
\allowdisplaybreaks
\begin{multline*}
\sum_{q\,\geq\,3}q!\;\sum_{i_{1}+i_{2}+i_{3}=q}\;\sum_{j_{1}+j_{2}+j_{3}=q}\frac{\left|\alpha _{i_{1},i_{2}}\beta _{i_{3}}(u)\right|}{i_{1}!i_{2}!i_{3}!}\frac{\left|\alpha _{j_{1},j_{2}}\beta _{j_{3}}(u)\right|}{j_{1}!j_{2}!j_{3}!}\varepsilon^{q} \\
\begin{aligned}
&=\sum_{q\,\geq\,3}\varepsilon^{q}\;\sum_{i_{1}+i_{2}+i_{3}=q}\;\sum_{j_{1}+j_{2}+j_{3}=q}\sqrt{\left(i_{1}+i_{2}+i_{3}\right)!}\sqrt{\left(j_{1}+j_{2}+j_{3}\right)!} \\
&\qquad\qquad\qquad\times \frac{\left|\alpha_{i_{1},i_{2}}\beta_{i_{3}}(u)\right|}{i_{1}!i_{2}!i_{3}!}\frac{\left|\alpha_{j_{1},j_{2}}\beta_{j_{3}}(u)\right|}{j_{1}!j_{2}!j_{3}!} \\
&=\sum_{q\,\geq\,3}\varepsilon^{q}\;\sum_{i_{1}+i_{2}+i_{3}=q}\;\sum_{j_{1}+j_{2}+j_{3}=q}\sqrt{\frac{\left(i_{1}+i_{2}+i_{3}\right)!}{i_{1}!i_{2}!i_{3}!}}\sqrt{\frac{\left(j_{1}+j_{2}+j_{3}\right)!}{j_{1}!j_{2}!j_{3}!}} \\
&\qquad\qquad\qquad\times\frac{\left|\alpha _{i_{1},i_{2}}\beta _{i_{3}}(u)\right|}{\sqrt{i_{1}!i_{2}!i_{3}!}}\frac{\left|\alpha _{j_{1},j_{2}}\beta _{j_{3}}(u)\right|}{
\sqrt{j_{1}!j_{2}!j_{3}!}} \\
&\leq \sum_{q\,\geq\,3}\varepsilon^{q}\sqrt{\sum_{i_{1}+i_{2}+i_{3}=q}
\sum_{j_{1}+j_{2}+j_{3}=q}\frac{\left(i_{1}+i_{2}+i_{3}\right)!}{i_{1}!i_{2}!i_{3}!}
\frac{\left(j_{1}+j_{2}+j_{3}\right)!}{j_{1}!j_{2}!j_{3}!}} \\
&\quad \times \sqrt{\sum_{i_{1}+i_{2}+i_{3}=q}\;\sum_{j_{1}+j_{2}+j_{3}=q}
\frac{\left|\alpha _{i_{1},i_{2}}\beta _{i_{3}}(u)\right|^{2}}{i_{1}!i_{2}!i_{3}!}\frac{\left|\alpha _{j_{1},j_{2}}\beta_{j_{3}}(u)\right|^{2}}{j_{1}!j_{2}!j_{3}!}} \\
&\leq \sum_{q\,\geq\,3}K_{1}(q)\varepsilon ^{q}\sqrt{
\sum_{i_{1}+i_{2}+i_{3}=q}\;\sum_{j_{1}+j_{2}+j_{3}=q}\frac{\left(i_{1}+i_{2}+i_{3}\right)!}{i_{1}!i_{2}!i_{3}!}\frac{\left(j_{1}+j_{2}+j_{3}\right)!}{j_{1}!j_{2}!j_{3}!}}
\end{aligned}
\end{multline*}
\endgroup
where
\begin{equation*}
K_{1}(q)=\sum_{i_{1}+i_{2}+i_{3}=q}\frac{\left|\alpha _{i_{1},i_{2}}\beta _{i_{3}}(u)\right|^{2}}{i_{1}!i_{2}!i_{3}!}\leq const,\text{ }
\end{equation*}
uniformly in $q$ note that this is the variance of the $q-$th order chaos for the expansion of the boundary length for a unit variance spherical random field. On the other hand,
\begin{align*}
\sqrt{\sum_{i_{1}+i_{2}+i_{3}=q}\sum_{j_{1}+j_{2}+j_{3}=q}\frac{\left(i_{1}+i_{2}+i_{3}\right)!}{i_{1}!i_{2}!i_{3}!}\frac{\left(j_{1}+j_{2}+j_{3}\right)!}{j_{1}!j_{2}!j_{3}!}} 
=\sum_{i_{1}+i_{2}+i_{3}=q}\frac{\left(i_{1}+i_{2}+i_{3}\right)!}{i_{1}!i_{2}!i_{3}!}= 3^{q},
\end{align*}
whence
\begin{align*}
\sum_{q\,\geq\,3}K_{1}(q)\varepsilon^{q}\sqrt{\sum_{i_{1}+i_{2}+i_{3}=q}
\sum_{j_{1}+j_{2}+j_{3}=q}\frac{\left(i_{1}+i_{2}+i_{3}\right)!}{i_{1}!i_{2}!i_{3}!}
\frac{\left(j_{1}+j_{2}+j_{3}\right)!}{j_{1}!j_{2}!j_{3}!}} 
=\sum_{q\,\geq\,3}K_{1}(q)\varepsilon^{q}3^{q}<\infty \text{,}
\end{align*}
since one can choose $\varepsilon <\frac{1}{3}$. Consequently, we just proved that
\begin{equation}\label{var-high}
\sum_{q\,\geq\,3}\Var(\clC_{T}(u)[q])=O\left(T^{2-\frac52
\beta _{\ell ^{\star }}}\right)=o\left(T^{2-2 \beta _{\ell ^{\star }}}\right)
\end{equation}
thus concluding the proof of Proposition~\ref{var-higher-chaoses}.
\end{proof}

\subsection{Proof of Theorem~\ref{main-theorem}}

We will need the following well known result.

\begin{theo}[{\cite{DM79, Ta:79}}]\label{DMtheorem}
Let $\xi(t)$, $t\in\bbR$, be a real measurable mean-square continuous stationary Gaussian process with mean $\bbE
\left[\xi(t)\right]$ and covariance function $\rho(t-s)=\rho(|t-s|)=
\mathop{\rm Cov}(\xi(t),\xi(s))$. Moreover, assume that
\begin{equation}\label{eq:DM79}
\rho(t-s)=\frac{L(|t-s|)}{|t-s|^{\beta}}\,, \qquad \text{with} \quad 0<\beta<1,
\end{equation}
where $L$ is a slowly varying function. Let $F: \bbR \rightarrow
\bbR$ be a Borel function such that $\bbE\left[F(N)^2\right] <+\infty$, where $N$ is a standard Gaussian random variable. Then it is a well known fact that can be expanded as follows
\begin{equation*}
F(\xi)=\sum_{k=0}^\infty \frac{b_k}{k!}H_k(\xi)\,, \quad \text{where} \quad b_k=\int_{\bbR}F(\xi)H_k(\xi)\phi(\xi)d\xi\,.
\end{equation*}
Assume there exists an integer $r$, the so-called Hermitian rank, such that $ b_0=b_1=\cdots=b_{r-1}=0$ and $b_r\ne0$. Then, if $\beta\in(0,1/r)$, we have that the finite-dimensional distributions of the random process
\begin{equation*}
X_T(s)=\frac{1}{T^{1-\beta r /2}L(T)^{r/2}}\int_{0}^{Ts}\,\left[F(\xi(t))-b_0
\right] \, dt \,, \qquad 0\le s \le1\,,
\end{equation*}
converge weakly, as $T\rightarrow \infty$, to the ones of the Rosenblatt process of order $r$, that is
\begin{equation*}
X_{\beta}(s):=\frac{b_r}{r!}\int_{(\bbR^r)^{\prime}}\,\frac{ e^{i(\lambda_1+\cdots+\lambda_r)s}-1}{i(\lambda_1+\cdots+\lambda_r)}\frac{ W(d\lambda_1)\cdots W(d\lambda_r)}{|\lambda_1\cdots\lambda_r|^{(1-\beta)/2}}
\, dt \,, \qquad 0\le s \le1\,,
\end{equation*}
where $W$ is a complex Gaussian white noise.
\end{theo}

\begin{proof}[Proof of Theorem~\ref{main-theorem}]
Recall that $2\beta_{\ell^\star}<\min(\beta_0,1)$. From Lemma~\ref{lemvar1} we have
\[
\lim_{T\,\to\,\infty} \frac{\Var(\clC_T(u)[1])}{T^{2-2\beta_{\ell^\star}}} = 0\,.
\]
Moreover, thanks to Proposition~\ref{var-higher-chaoses},
\[
\lim_{T\,\to\,\infty} \frac{\sum_{q\,\ge\,3} \Var(\clC_T(u)[q])}{T^{2-2\beta_{\ell^\star}}} = 0\,,
\]
so that, recalling also Proposition~\ref{prop:var-2nd-chaos},
\begin{equation}\label{2chaosA}
\frac{\clC_T(u)}{T^{1-\beta_{\ell^\star}}}=\frac{\clC_T(u)[2]}{T^{1-\beta_{\ell^\star}}}+o_{\bbP}(1)\,.
\end{equation}
Moreover, since in $L^2(\Omega)$ we have the following equality (recall Remark~\ref{2chaos-hermite})
\begin{multline*}
\clC_T(u)[2]=\\
\frac{\sigma _{1}}{2}\sqrt{\frac{\pi }{2}}
\phi (u)\sum_{\ell} \frac{C_{\ell }(0)(2\ell +1)}{4\pi } \left\{\left(u^{2}-1\right)+\frac{\lambda _{\ell }/2}{
\sigma_{1}^{2}}\right\} \int_0^T\int_{\bbS^{2}}H_{2}\left(\widehat Z_{\ell}(x,t)\right)dxdt\text{,}
\end{multline*}
it holds that
\begin{multline} \label{eq:var2spectral}
\frac{\clC_T(u)[2]}{T^{1-\beta_{\ell^\star}}}=
\sum_{\ell\,\in\,\clI^\star} \frac{\sigma_1}{2}\sqrt{\frac\pi2} \phi(u) C_\ell(0) \left\{\left(u^{2}-1\right)+\frac{\lambda_{\ell}/2}{\sigma_{1}^{2}}\right\}\\
\times \sum_{m=-\ell}^{\ell} \frac{1}{T^{1-\beta_{\ell^\star}}}\int_{0}^T H_2\left(\widehat a_{\ell m}(t)\right)dt +o_{\bbP}(1)\,,
\end{multline}
where $\hat a_{\ell m}(t):= a_{l m}(t) / \sqrt{C_\ell(0)}$. Indeed, recalling Proposition~\ref{prop:var-2nd-chaos}, we have that
\begin{equation}\label{th21}
\lim_{T\,\rightarrow\,\infty }\frac{\Var\left(\clC_{T}(u)[2]\right) }{T^{2-2\beta_{\ell^{\star}}}} =\frac{
\sigma_{1}^{2}\pi }{4}\phi^{2}(u)\sum_{\ell\,\in\,\clI^\star}\frac{(2\ell+1)^{2}C_\ell(0)^2}{(1-2\beta_\ell)(1-\beta_\ell)}
\left\{\left(u^{2}-1\right)+\frac{\lambda _{\ell}/2}{\sigma_{1}^{2}}\right\}^2\,.
\end{equation}
and hence that
\begin{multline*}
\lim_{T\,\to\,\infty}\mathbb E\left[\left(\frac{\clC_T(u)[2]}{T^{1-\beta_{\ell^\star}}} -\frac{1}{T^{1-\beta_{\ell^\star}}}\sum_{\ell\,\in\,\clI^\star}\sum_{m=-\ell}^\ell \frac{J_2(u)}{2} C_\ell(0)\int_0^T H_2(\hat a_{\ell, m}(t))\, dt \right)^2\right]\\
=\lim_{T\,\to\,\infty}\frac{\sigma_{1}^{2}\pi }{4\,T^{2-2\beta_{\ell^\star}}}\phi ^{2}(u)\sum_{\ell\,\notin\,\clI^\star}(2\ell +1)^{2}\left\{\left(u^{2}-1\right)+\frac{\lambda_{\ell }/2}{\sigma_{1}^{2}}\right\} ^{2}\int_{[0,T]^{2}}C_{\ell }(t-s)^{2}dtds\,.
\end{multline*}
From~\eqref{2chaosA} and~\eqref{eq:var2spectral}, in order to understand the asymptotic distribution of $\clC_T(u)$, it suffices to investigate the leading term on the right hand side of~\eqref{eq:var2spectral}. Recall Assumption~\ref{basic3}, for $\ell\in \clI^\star$ we have that
\[
C_{\ell}(\tau)= \frac{G_\ell(\tau)}{(1+|\tau|)^{\beta_{\ell^\star}}} \,,
\]
where in particular $G_\ell$ is a slowly varying function. Hence, setting $\xi(t)=a_{\ell, m}(t)$ in Theorem~\ref{DMtheorem}, we automatically have that $\rho=\rho_\ell=C_{\ell}$, $L=L_\ell=G_\ell$ and, as a consequence, that
\begin{equation*}
X_T^{\ell, m}:=\frac{1}{T^{1-\beta_{\ell^\star}}}\int_0^T H_2(\widehat a_{\ell, m}(t))\, dt
\overset{d}{\longrightarrow} \frac{X_{m;\beta_{\ell^\star}}}{a\left(\beta_{\ell^\star}\right)}\,, \qquad \text{as } T \rightarrow
\infty\,,
\end{equation*}
for all $m=-\ell,\,\dots,\,\ell$, where, for each $m$, $X_{m;\beta_{\ell^\star}}$ is a standard Rosenblatt random variable~\eqref{Xbeta} of parameter $\beta_{\ell^\star}$. Moreover, since the $X_T^{\ell, m}$ are all independent for each $T$ we have that

\begin{align*}
\widetilde{\clC}_{T}(u)=&\sqrt{\frac{T^{2-2\beta_{\ell^\star}}}{\Var \clC_{T}(u)[2]}} \sum_{\ell\,\in\,\clI^\star} \frac{\sigma_1}{2}\sqrt{\frac\pi2} \phi(u) C_\ell(0) \left\{\left(u^{2}-1\right)+\frac{\lambda _{\ell}/2}{\sigma_{1}^{2}}\right\}\\
&\qquad\times \sum_{m=-\ell}^{\ell} \frac{\int_{0}^T H_2\left(\widehat a_{lm}(t)\right)dt}{T^{1-\beta_{\ell^\star}}}+o_{\bbP}(1)
\\
\mathop{\longrightarrow}^d &\left(\frac{\sigma_{1}^{2}\pi }{4}\phi ^{2}(u)\sum_{\ell\, \in\,\clI^\star} \frac{(2\ell+1)^{2} C_\ell(0)^2}{(1-2\beta_\ell)(1-\beta_\ell)} \left\{ \left(u^{2}-1\right)+\frac{\lambda _{\ell}/2}{\sigma_{1}^{2}}\right\}^2 \right)^{-1/2}
\\
\qquad\times\,&\sum_{\ell\,\in\,\clI^\star} \frac{\sigma_1^2}{2}\sqrt{\frac\pi2} \phi(u) C_\ell(0) \left\{\left(u^{2}-1\right)+\frac{\lambda _{\ell}/2}{\sigma_{1}^{2}}\right\}\sum_{m=-\ell}^{\ell} \frac{X_{m,\beta_{\ell^\star}}}{a\left(\beta_{\ell^\star}\right)}
\\
\mathop{=}^{d}\,&\sum_{\ell \in \clI^\star} \frac{C_\ell(0)}{\sqrt{v^\star}}\left\{\left(u^{2}-1\right)+\frac{\lambda _{\ell}/2}{\sigma_{1}^{2}}\right\}V_{2\ell+1}(1,\,\dots,\,1\,;\,\beta_{\ell^\star})\,,
\\
\intertext{where}
v^\star\,=\,& a\left(\beta_{\ell^\star}\right)^2 \sum_{\ell\, \in\,\clI^\star} \frac{2(2\ell+1)^{2} C_\ell(0)^2}{(1-2\beta_\ell)(1-\beta_\ell)} \left\{ \left(u^{2}-1\right)+\frac{\lambda _{\ell}/2}{\sigma_{1}^{2}}\right\}^2
\end{align*}
and the proof of Theorem~\ref{main-theorem} is concluded.
\end{proof}

Some auxiliary results are collected in the four appendixes that follow.


\appendix

\section{Covariance structure}
\label{app-cov}

In this Section we collect technical results on the covariance structure of the field $(Z, \nabla Z)$.


\begin{lemm}\label{lemcov}
Let $Z$ be a space-time spherical random field satisfying Assumption~\ref{basic} and Assumption~\ref{regularT}. Then for all points $x=(\theta _{x},\varphi _{x}),$ $y=(\theta _{y},\varphi _{y})\in \bbS^{2}\setminus \{N,S\}$, the covariance structure of $(Z, \nabla Z)$ is
\[
\bbE\left[Z(x,t)Z(y,s)\right] = \Gamma(\langle x, y\rangle, t-s),
\]
\begin{multline*}
\bbE\left[Z(x,t)\partial _{1;y}Z(y,s)\right]
\\
=\sum_{\ell =0}^{\infty }\frac{2\ell +1}{4\pi }C_{\ell }(t-s)P_{\ell }^{\prime }(\left\langle x,y\right\rangle)\left\{ -\cos \theta _{x}\sin
\theta _{y}+\sin \theta _{x}\cos \theta _{y}\cos (\varphi _{x}-\varphi _{y})\right\} \text{,}
\end{multline*}
\begin{equation*}
\bbE\left[Z(x,t)\partial _{2;y}Z(y,s)\right]
=\sum_{\ell =0}^{\infty }\frac{2\ell +1}{4\pi }C_{\ell }(t-s)P_{\ell }^{\prime }(\left\langle x,y\right\rangle)\left\{ \sin \theta _{x}\sin (\varphi _{x}-\varphi _{y})\right\} \text{,}
\end{equation*}
and moreover

\begin{multline*}
\bbE\left[\partial _{1;x}Z(x,t)\partial _{1;y}Z(y,s)\right] =
\\
\begin{aligned}
&\sum_{\ell =0}^{\infty }\frac{2\ell +1}{4\pi }C_{\ell }(t-s)P_{\ell }^{\prime \prime }(\left\langle x,y\right\rangle)\left\{ -\cos \theta _{x}\sin \theta _{y}+\sin \theta _{x}\cos \theta _{y}\cos (\varphi _{x}-\varphi _{y})\right\} \\
&\times \left\{ -\sin \theta _{x}\cos \theta _{y}+\cos \theta _{x}\sin
\theta _{y}\cos (\varphi _{x}-\varphi _{y})\right\} \\
&+\sum_{\ell =0}^{\infty }\frac{2\ell +1}{4\pi }C_{\ell }(t-s)P_{\ell }^{\prime }(\left\langle x,y\right\rangle)\left\{ \sin \theta _{x}\sin
\theta _{y}+\cos \theta _{x}\cos \theta _{y}\cos (\varphi _{x}-\varphi _{y})\right\} \text{,}
\end{aligned}
\end{multline*}
\begin{multline*}
\bbE\left[\partial _{1;x}Z(x,t)\partial _{2;y}Z(y,s)\right] =
\\
\sum_{\ell =0}^{\infty }\frac{2\ell +1}{4\pi }C_{\ell }(t-s)P_{\ell }^{\prime \prime }(\left\langle x,y\right\rangle) \\
\times \left\{ -\sin \theta _{x}\cos \theta _{y}+\cos \theta _{x}\sin
\theta _{y}\cos (\varphi _{x}-\varphi _{y})\right\} \left\{ \sin \theta _{x}\sin \theta _{y}\sin (\varphi _{x}-\varphi _{y})\right\} \\
+\sum_{\ell =0}^{\infty }\frac{2\ell +1}{4\pi }C_{\ell }(t-s)P_{\ell }^{\prime }(\left\langle x,y\right\rangle)\cos \theta _{x}\sin \theta _{y}\sin (\varphi _{x}-\varphi _{y})\text{ },
\end{multline*}
\begin{multline*}
\bbE\left[\partial _{2;x}Z(x,t)\partial _{2;y}Z(y,s)\right] =
\\
-\sum_{\ell =0}^{\infty }\frac{2\ell +1}{4\pi }C_{\ell }(t-s)P_{\ell }^{\prime \prime }(\left\langle x,y\right\rangle)\sin \theta _{x}\sin
\theta _{y}\sin ^{2}(\varphi _{x}-\varphi _{y}) \\
+\sum_{\ell =0}^{\infty }\frac{2\ell +1}{4\pi }C_{\ell }(t-s)P_{\ell }^{\prime }(\left\langle x,y\right\rangle)\sin \theta _{x}\sin \theta _{y}\cos (\varphi _{x}-\varphi _{y})\text{ }.
\end{multline*}
\end{lemm}

The proof of these results is entirely analogous to the one given in~\cite[App\-endix]{CM20} and hence omitted; note that in the latter reference the definition of $\partial _{2;y}$ differs by a factor $\frac{1}{\sin \theta_{y}}$, i.e., covariant derivatives are used in the computations.

\section{Measurability issues}

Let us recall that the space-time spherical random field $Z$ satisfies Assumptions~\ref{basic} and~\ref{regularT}.

Let $u\in \bbR$ be a fixed threshold, for $ t\in \bbR$ we consider the level set $Z(\cdot, t)^{-1}(u):=\{ x\in \bbS^2 : Z(x,t) = u\}$ which is an a.s. $\clC^1$ manifold of dimension~$1$. Indeed, for every $t$, $Z(\cdot, t)\in \clC^1(\bbS^2)$, and for every $x\in Z(\cdot, t)^{-1}(u)$ the covariance of $(Z(x,t), \nabla_x Z(x,t))$ is non-degenerate, hence Bulinskaya's lemma/Ylvisaker's lemma ensures that there exists $
\Omega_t\subseteq \Omega$, $\bbP(\Omega_t)=1$, such that for every $
\omega\in \Omega_t$, the value $u$ is regular for $Z(\cdot, t)(\omega)$, i.e.
\begin{equation*}
\nabla_x Z(x, t)(\omega)\ne 0 \text{ for every } x \text{ such that } Z(x,t)(\omega)=u.
\end{equation*}
Hence, on $\Omega_t$ we can define
\begin{equation*}
\clL_u(t):=\clH^1\left(Z(\cdot, t)^{-1}(u)\right).
\end{equation*}
\begin{rema}
Clearly, the previous discussion does \emph{not} ensure that there exists $\Omega^0\subseteq \Omega$, $\bbP(\Omega^0)=1$ s.t., for every $\omega\in \Omega^0$ the value $u$ is regular for $Z(\cdot, t)(\omega)$ for every $t$. In other words, we \emph{cannot} guarantee that a boundary length process $\clL_t(u)$, indexed by $t\in \bbR$, exists. However, in Lemmas~\ref{lemmeas} and~\ref{lemreg} to follow, we are going to prove that, for any $T>0$, the value $u$ is regular for $Z(\cdot, t)(\omega)$ outside of a negligible subset of $(\omega, t)\in \Omega\times [0,T]$. This is sufficient to define our quantity of interest $\clC_T(u)$ for every $T>0$ in~\eqref{rap_int}, see also Lemma~\ref{lem_basic}.
\end{rema}

By stationarity, the law of $\clL_u(t)$ does not depend on $t$, in particular $\bbE[\clL_u(t)]$ does not depend on $t$, and can be computed via the Kac--Rice formula~\cite[Theorem~6.8]{AW:09} or the Gaussian Kinematic Formula~\cite[Theorem~13.2.1]{adlertaylor} to be
\begin{equation*}
\bbE[\clL_u(t)] = \sigma_1 \cdot 2\pi e^{-u^2/2},
\end{equation*}
where $\sigma_1$ is defined as in~\eqref{sigma1}. In order to define our functional of interest and prove that it is indeed a random variable, we need the following technical result.

\begin{lemm}\label{lemmeas}
For every $T>0$, the set
\begin{align*}
A^T {:=}&\, \left \{ (\omega, t)\in \Omega\times [0,T]: \{ x\in
\bbS^2 : |Z(x,t)(\omega) - u | =0, |\nabla Z(x,t)(\omega) |=0\} =\emptyset \right \} \\
=&\,\left\{ (\omega, t)\in \Omega \times [0,T] : \text{ the value } u \text{ is regular for } Z(\cdot, t)(\omega)\right\}.
\end{align*}
is measurable, i.e., $A^T\in \clF\otimes \clB([0,T])$.
\end{lemm}

\begin{proof}
Fix $\{x_i\}_{i\,\in\,\bbN}$ to be a \emph{dense} sequence in $\bbS^2$; for $n,k\in \bbN$, define the set
%{\small
\begin{multline}
A^T_{n,k}:= \bigcap_{i\,\in\,\bbN} \left\{ (\omega, t)\in \Omega\times [0,T]: |Z(x_i, t, \omega) - u | \ge \frac{1}{n} \right \} \\
\cup \left \{(\omega, t)\in \Omega\times [0,T] : |\nabla Z(x_i, t, \omega)| \ge \frac{1}{k} \right \}
\end{multline}
%}
which is measurable by construction, i.e. $A^T_{n,k}\in \mathcal F\otimes \mathcal B([0,T])$. We will show that $A^T$ equals
\begin{equation}
\bigcup_{k,\,n\,\in\,\bbN} A^T_{n,k} =: \widehat A^T,
\end{equation}
so that, in particular, $A^T$ is measurable. Indeed,
\begin{equation}
A^T\subseteq \widehat A^T,
\end{equation}
because if $(\omega,t)\in A^T$, then
\begin{equation}
\inf_{x\,\in\,\bbS^2\,:\,Z(x,t)(\omega)=u} |\nabla Z(x,t)(\omega)| =: l >0
\end{equation}
hence by continuity of $x\mapsto Z(x,t)(\omega)$ and $x\mapsto \nabla Z(x,t)(\omega)$ there exist $\tilde k, \tilde n\in \bbN$ s.t.
\begin{equation}
\inf_{x\,\in\,\bbS^2\,:\,|Z(x,t)(\omega) - u| < \frac{1}{\tilde n}} |\nabla Z(x,t)(\omega)| > l/2 > \frac{1}{\tilde k}.
\end{equation}
Thus $(\omega,t)\in A^T_{\tilde n, \tilde k}$. On the other hand, if $(\omega,t)\in \widehat A^T$, then there exist $\tilde n, \tilde k\in \bbN$ s.t. $(\omega,t)\in A^T_{\tilde n, \tilde k}$, that is,
\begin{equation}
\inf_{i\,\in\,\bbN\,:\,|Z(x_i, t)(\omega) - u|\,<\,1/\tilde n} |\nabla Z(x_i, t)(\omega) | \ge 1/\tilde k.
\end{equation}
By continuity of $Z(\cdot,t)(\omega)$ and $\nabla Z(\cdot, t)(\omega)$, and density of the sequence $\{ x_i\}_i$,
\begin{equation}
\inf_{x\,\in\,\bbS^2\,:\,|Z(x, t)(\omega) - u|< 1/2\tilde n} |\nabla Z(x, t)(\omega) | \ge 1/\tilde k
\end{equation}
hence
\begin{equation}
\inf_{x\,\in\,\bbS^2\,:\,Z(x, t)(\omega) = u} |\nabla Z(x, t)(\omega) | \ge 1/\tilde k >0
\end{equation}
and $(\omega,t)\in A^T$.
\end{proof}

\begin{lemm}\label{lemreg}
Let $T>0$. There exists $\tilde \Omega_T \subseteq \Omega$, $\bbP(\tilde \Omega_T)=1$, such that for every $\omega\in \tilde\Omega_T$ there exists $I_T(\omega)\subseteq [0,T]$, $\mathrm{Leb} (I_T(\omega))=T$, such that the value $u$ is regular for $Z(\cdot, t)(\omega) $ for every $t\in I_T(\omega)$.
\end{lemm}

We are now in a position to prove Lemma~\ref{lem_basic}. (Lemma~\ref{lemreg} will be established right after.)

\begin{proof}[Proof of Lemma~\ref{lem_basic}]
In view of Lemma~\ref{lemreg}, let us define $\tilde \Omega:=\cap_{n\,\in\,\bbN} \tilde \Omega_n$ ($\bbP(\tilde \Omega)=1$), hence for every $\omega\in \tilde \Omega$ there exists $I(\omega)\subseteq [0,+\infty)$ ($\mathrm{Leb}(I(\omega)^c) = 0$) such that the value $u$ is regular for $ Z(\cdot, t)(\omega)$ for every $t\in I(\omega)$.
\end{proof}

On $\tilde \Omega$ we can define the quantity
\begin{equation*}
\clC_{T}(u)(\omega):=\,\int_{0}^{T}\,\Big(\clL_{u}(t)(\omega)-
\bbE[\clL_{u}(t)]\Big)\,dt
\end{equation*}
for every $T>0$, which is a random variable.

\begin{proof}[Proof of Lemma~\ref{lemreg}]
For $T>0$, we consider the measure space

$([0,T], \mathcal B([0,T]), \text{Leb}_{[0,T]})$, and define
\[
A^T := \{ (\omega, t)\in \Omega \times [0,T] : \text{ the value } u \text{ is regular for } Z(\cdot, t)(\omega)\}.
\]
From Lemma~\ref{lemmeas}, $A^T$ is measurable ($A^T \in \mathfrak F\otimes \mathcal B([0,T])$). In particular, the section
\[
A^T_t = \{ \omega\in \Omega : \text{ the value } u \text{ is regular for } Z(\cdot, t)(\omega)\}
\]
is a measurable set and contains $\Omega_t$ ensuring that $\bbP(A^T_t)=1$. A standard application of Fubini's theorem gives
\begin{equation}\label{meas0}
T = \int_0^T \bbP(A^T_t)\,dt = \mathbb E \left [\int_0^T 1_{A^T}(\omega,t)\,dt
\right]
\end{equation}
implying that there exists $\tilde \Omega_T\subseteq \Omega$, $\bbP(\tilde \Omega_T) =1$, such that for every $\omega\in \tilde \Omega_T$ we have $
\int_0^T 1_{A^T}(\omega,t)\,dt = T. $
\end{proof}

\section{Square integrability}
\label{app-square}

In this Section first we prove that $\clC_T(u)$ is square integrable. By a standard application of Jensen's inequality and the stationarity of the model we have
\begin{equation}
\bbE\left[\clC_T(u)^2\right] \le T^2 \rVar(\clL_u(0))
\end{equation}
for any $T>0$. Hence it suffices to prove that $\clL_u(0)$ is square integrable (clearly, it is equivalent to show that $\clL_u(t)$ is so, for any $t\in\bbR$).

Recall the definition of $\epsilon$-approximating random variables $\clL_u^\epsilon(t)$ in~\eqref{Leps}.

\begin{lemm}\label{approx2}
Let $t\in \bbR$ be fixed. As $\epsilon\to 0$,
\begin{equation}\label{convergence}
\clL_u^\epsilon(t) \to \text{length}\left(Z(\cdot, t)^{-1}(u)\right)=\clL _u(t)
\end{equation}
both a.s. and in $L^2(\bbP)$.
\end{lemm}

\begin{proof}
The following conditions are satisfied:
\begin{enumerate}
\item \label{punto1}
for every fixed $t\in \bbR$, the random field $Z(\cdot,t)$ is with probability one a Morse function on $\bbS^2$, see Section~\ref{subsec-object} for details;
\item \label{punto2}
the covariance function $\Gamma$ of the field is at least twice continuously differentiable with strictly positive second-order derivative in a neighborhood of the origin, meaning that
\begin{equation*}
\sum_{\ell }\ell ^{2}\frac{2\ell +1}{4\pi }C_{\ell }(0)<\infty \text{.}
\end{equation*}
Note that this expression is strictly positive unless $C_{\ell }(0)=0$ for all $
\ell \geq 1.$
\end{enumerate}

Also, \eqref{punto2} implies that the second derivative of the covariance function is continuous at the origin, and hence
\begin{equation*}
1-\Gamma (\cos \theta,0)=\Gamma ^{\prime \prime }(0,0)\theta ^{2}+o(\theta ^{2})\text{, as }\theta \rightarrow 0\text{.}
\end{equation*}
We recall incidentally that
\begin{equation*}
\frac{\partial ^{2}}{\partial \theta ^{2}}\Gamma (\cos \theta)=\sum_{\ell }
\frac{2\ell +1}{4\pi }C_{\ell }P_{\ell }^{\prime \prime }(\cos \theta)\sin ^{2}\theta -\sum_{\ell }\frac{2\ell +1}{4\pi }C_{\ell }P_{\ell }^{\prime }(\cos \theta)\cos \theta
\end{equation*}
which by Cauchy--Schwartz inequality has a unique maximum for $\theta =0,$ given by
\begin{equation*}
\Gamma ^{\prime \prime }(0,0)=\sum_{\ell }\frac{2\ell +1}{4\pi }\frac{
\lambda _{\ell }}{2}C_{\ell }\text{.}
\end{equation*}
For notational simplicity we prove that the $L^{2}$-expansion holds at $u=0;$ the proof for different values is identical. Our argument is quite standard, see for instance~\cite{MRW}.

We know that the boundary length is defined almost-surely by
\begin{align*}
\clL_{0}(t) &\mathcal{=}\lim_{\varepsilon\,\rightarrow\,0}\clL _{0\,;\,\varepsilon }(t)\text{, } \\
\clL_{0\,;\,\varepsilon }(t) &:=\int_{\bbS^{2}}\delta _{\varepsilon }(Z(x,t))\left\Vert \nabla Z(x,t)\right\Vert dx\text{,}
\end{align*}
where
\begin{equation*}
\delta _{\varepsilon }(Z(x,t)):=\left\{
\begin{array}{c}
0\text{ for }x:Z(x,t)>\varepsilon \\
\frac{1}{2\varepsilon }\text{ for }x:Z(x,t)\leq \varepsilon
\end{array}
\right.
\end{equation*}
and the almost-sure convergence follows from the standard arguments \cite[Lem\-ma~3.1]{RW08}. Indeed, because $\delta _{\varepsilon }$ is integrable and $Z(\cdot)$ is Morse we have, using the coarea formula for a fixed $t\in
\bbR$ (see i.e., \cite[p.~169]{adlertaylor})
\begin{equation*}
\int_{\bbS^{2}}\delta_{\varepsilon }(Z(x,t))\left\Vert \nabla Z(x,t)\right\Vert dx=\int_{\bbR}\left\{ \int_{Z^{-1}(s,t)}\delta_{\varepsilon }(Z(x,t))dx\right\} ds
\end{equation*}
and thus we obtain
\begin{equation*}
\int_{\bbR}\left\{ \int_{Z^{-1}(s,t)}\delta_{\varepsilon }(Z(x,t))dx\right\} ds=\frac{1}{2\varepsilon }\int_{-\varepsilon }^{\varepsilon }\textrm{length}\left[Z^{-1}(s,t)\right] ds\rightarrow \textrm{length}\left[Z^{-1}(0,t)\right],
\end{equation*}
as $\varepsilon \rightarrow 0$, because the function $s\rightarrow \textrm{length}[Z^{-1}(s,t)] $ is continuous for Morse functions, see~\ref{punto1}. In particular, \eqref{convergence} holds a.s.

We now want to show that the convergence occurs also in the $L^{2}$ sense; because convergence holds almost surely, it is enough to show that
\begin{equation*}
\lim_{\varepsilon\,\rightarrow\,0}\bbE\left[\clL_{0;\varepsilon }^{2}(t)\right] =\bbE\left[\clL_{0}^{2}(t)\right] \text{.}
\end{equation*}
Indeed, we have that
\begin{align*}
\lim_{\varepsilon\,\rightarrow\,0}\bbE\left[\left(\clL_{0}(t)-\clL_{0;\varepsilon }(t)\right)^{2}\right] &=\lim_{\varepsilon\,
\rightarrow\,0}\bbE\left[\left(\clL_{0}^{2}(t)+\clL_{0;\varepsilon}^{2}(t)-2\clL_{0}(t)\clL_{0\,;\,\varepsilon}(t)
\right)\right] \\
&=2\bbE\left[\clL_{0}^{2}(t)\right] -2\lim_{\varepsilon
\,\rightarrow\,0}\bbE\left[\clL_{0}(t)\clL_{0\,;\,\varepsilon }(t)\right] =0\text{,}
\end{align*}
because by Fatou's Lemma and Cauchy--Schwartz inequality
\begin{equation*}
\bbE\left[\clL_{0}^{2}(t)\right] \leq \lim_{\varepsilon
\,\rightarrow\,0}\inf \bbE\left[\clL_{0}(t)\clL _{0\,;\,\varepsilon}(t)\right] \leq \lim_{\varepsilon\,\rightarrow\,0}\sqrt{
\bbE\left[\clL_{0}^{2}(t)\right] \bbE\left[\clL _{0\,;\,\varepsilon}^{2}(t)\right] }=\bbE\left[\clL_{0}^{2}(t)
\right] \text{.}
\end{equation*}
Note that, by the coarea formula
\begin{align*}
\bbE\left[\clL_{0\,;\,\varepsilon}^{2}(t)\right] &=\bbE
\left[\left\{\int_{\bbS^{2}}\left\{\delta_{\varepsilon }(Z(x,t))\left\Vert \nabla Z(x,t)\right\Vert \right\} dx\right\} ^{2}\right]
\\
&=\bbE\left[\left\{\int_{\bbR}\int_{Z(x,t)=u}\delta _{\varepsilon }(Z(x,t))dxdu\right\}^{2}\right] \\
&=\bbE\left[\left\{\int_{\bbR}\clL_{u}(t)\delta _{\varepsilon }(u)du\right\}^{2}\right] \text{,}
\end{align*}
where as before by $\clL_{u}(t)$ we denote the length of the set $ Z(x,t)=u.$ We can now show that the application $u\rightarrow \bbE
\left[\clL_{u}^{2}(t)\right] $, or more explicitly
\begin{align*}
\bbE\left[\clL_{u}^{2}(t)\right] &=\int_{\bbS^{2}\times
\bbS^{2}}\bbE\Big[ \left\Vert \nabla Z(x_{1},t)\right\Vert \left\Vert \nabla Z(x_{2},t)\right\Vert \,\Big|\, Z(x_{1},t)=u,Z(x_{2},t)=u\Big]\\
&\quad\times \phi _{Z(x_{1},t),Z(x_{2},t)}(u,u)dx_{1}dx_{2}\\
&=8\pi ^{2}\int_{0}^{\pi }\text{ }\bbE\Big[ \left\Vert \nabla Z(N,t)\right\Vert \left\Vert \nabla Z(y(\theta),t)\right\Vert \,\Big|\, Z(N,t)=u,Z(y(\theta),t)=u\Big]\\
&\qquad\;\phi _{Z(N,t),Z(y(\theta),t)}(u,u)\sin
\theta d\theta \text{,}
\end{align*}
is continuous. The integrand function is obviously continuous in $u,$ and thus to check the latter statement it is enough to use Dominated Convergence Theorem. We first note that
\begin{multline*}
\phi _{Z(N,t),Z(y(\theta),t)}(u,u)\sin \theta\\
\leq \phi _{Z(N,t),Z(y(\theta),t)}(0,0)\sin \theta 
=\frac{1}{2\pi \sqrt{1-\Gamma^{2}(\cos \theta,0)}}\sin \theta =O(1)\text{,}
\end{multline*}
uniformly over $\theta $, because
\begin{equation*}
1-\Gamma ^{2}(\cos \theta,0)=(1+\Gamma (\cos \theta,0))(1-\Gamma (\cos
\theta,0))\geq \frac{c}{\theta ^{2}}\text{.}
\end{equation*}
On the other hand, to evaluate
\begin{equation*}
\bbE\Big[\left\Vert \nabla Z(x_{1},t)\right\Vert \left\Vert
\nabla Z(x_{2},t)\right\Vert \,|\, Z(Nt)=u,Z(y(\theta),t)=u\Big]
\end{equation*}
we can use Cauchy--Schwartz inequality, and bound
\begin{multline*}
\bbE\left[ w_{i}^{2}\,\middle|\, Z(N,t)=u,Z(y(\theta),t)=u
\right] \\
=\Var\big[w_{i}\,\big|\, Z(N,t)=u,Z(y(\theta),t)=u\big] +\big\{ \bbE\big[w_{i}\,\big|\, Z(N,t)=u,Z(y(\theta),t)=u
\big] \big\} ^{2},
\end{multline*}
for $i=1,2,3,4,$ where
\begin{equation*}
\left(
\begin{array}{c}
w_{1} \\
w_{2} \\
w_{3} \\
w_{4}
\end{array}
\right) :=\left(
\begin{array}{c}
\nabla Z(x_{1},t) \\
\nabla Z(x_{2}t)
\end{array}
\right) \text{.}
\end{equation*}
It is a standard fact for Gaussian conditional distributions that
\begin{equation*}
\Var\left[\left. w_{i}\right\vert Z(N,t)=u,Z(y(\theta),t)=u\right] \leq Var
\left[w_{i}\right] <\infty \text{.}
\end{equation*}
Similarly, standard results on Gaussian conditional expectations give (compare \cite[App\-endix A]{Wig},)
\begin{equation*}
\bbE\left[\left.
\begin{array}{c}
w_{1} \\
w_{2} \\
w_{3} \\
w_{4}
\end{array}
\right\vert Z(N,t)=u,Z(y(\theta)t)=u\right] =B_{\ell }^{T}(\theta)A_{\ell }^{-1}(\theta)\left(
\begin{array}{c}
u \\
u
\end{array}
\right) \text{,}
\end{equation*}
where
\begin{align*}
B_{\ell }^{T}(\theta) &=\left(
\begin{array}{cc}
\sum_{\ell }\frac{2\ell +1}{4\pi }C_{\ell }P_{\ell }^{\prime }(\cos \theta)\sin \theta & 0 \\
0 & 0 \\
0 & \sum_{\ell }\frac{2\ell +1}{4\pi }C_{\ell }P_{\ell }^{\prime }(\cos
\theta)\sin \theta \\
0 & 0
\end{array}
\right) \text{, } \\
A_{\ell }^{-1}(\theta) &=\frac{1}{1-\Gamma ^{2}(\cos \theta,0)}\left(
\begin{array}{cc}
1 & -\Gamma (\cos \theta,0) \\
-\Gamma (\cos \theta,0) & 1
\end{array}
\right) \text{,}
\end{align*}
That we obtain for the conditional expected value
%{\small
\begin{multline*}
\frac{1}{1-\Gamma ^{2}(\cos \theta,0)}\\
\left(
\begin{array}{cc}
-\sum_{\ell }\frac{2\ell +1}{4\pi }C_{\ell }P_{\ell }^{\prime }(\cos \theta)\sin \theta & \Gamma (\cos \theta)\sum_{\ell }\frac{2\ell +1}{4\pi } C_{\ell }P_{\ell }^{\prime }(\cos \theta)\sin \theta \\
0 & 0 \\
-\Gamma (\cos \theta)\sum_{\ell }\frac{2\ell +1}{4\pi }C_{\ell }P_{\ell }^{\prime }(\cos \theta)\sin \theta & -\sum_{\ell }\frac{2\ell +1}{4\pi } C_{\ell }P_{\ell }^{\prime }(\cos \theta)\sin \theta \\
0 & 0
\end{array}
\right) \left(
\begin{array}{c}
u \\
u
\end{array}
\right)
\\
\begin{aligned}
&=\frac{1}{1-\Gamma ^{2}(\cos \theta)}\left(
\begin{array}{c}
u(\Gamma (\cos \theta,0)-1)\sum_{\ell }\frac{2\ell +1}{4\pi }C_{\ell }P_{\ell }^{\prime }(\cos \theta)\sin \theta \\
0 \\
u(1-\Gamma (\cos \theta,0))\sum_{\ell }\frac{2\ell +1}{4\pi }C_{\ell }P_{\ell }^{\prime }(\cos \theta)\sin \theta \\
0
\end{array}
\right)
\\
&=\frac{1}{1+\Gamma (\cos \theta,0)}\left(
\begin{array}{c}
-u\sum_{\ell }\frac{2\ell +1}{4\pi }C_{\ell }P_{\ell }^{\prime }(\cos \theta)\sin \theta \\
0 \\
u\sum_{\ell }\frac{2\ell +1}{4\pi }C_{\ell }P_{\ell }^{\prime }(\cos \theta)\sin \theta \\
0
\end{array}
\right) \text{.}
\end{aligned}
\end{multline*}
This vector function is immediately seen to be uniformly bounded over $
\theta,$ whence the Dominated Convergence Theorem holds. To conclude the proof, we note that
\begin{align*}
\bbE\left[\clL_{0}^{2}(t)\right] &\leq \lim
\inf_{\varepsilon\,\rightarrow\,0}\bbE\left[\left\{ \int_{\bbS ^{2}}\left\{ \delta _{\varepsilon }(Z(x,t))\left\Vert \nabla Z(x,t)\right\Vert \right\} dx\right\} ^{2}\right] \\
&=\lim \inf_{\varepsilon\,\rightarrow\,0}\bbE\left[\clL _{0;\varepsilon }^{2}(t)\right] \leq \lim \sup_{\varepsilon \,\rightarrow\,0}
\bbE\left[\clL_{0;\varepsilon }^{2}(t)\right]
\end{align*}
(by Fatou's Lemma and definitions) and then
\begin{align*}
&=\lim \sup_{\varepsilon\,\rightarrow\,0}\bbE\left[\left\{ \int_{
\bbS^{2}}\left\{ \delta _{\varepsilon }(Z(x,t))\left\Vert \nabla Z(x,t)\right\Vert \right\} dx\right\} ^{2}\right] \\
&=\lim \sup_{\varepsilon\,\rightarrow\,0}\bbE\left[\left\{ \int_{\bbR}\clL_{u}(t)\delta_{\varepsilon }(u)du\right\} ^{2}\right]
\end{align*}
(by co-area formula) and
\begin{equation*}
\leq \lim \sup_{\varepsilon\,\rightarrow\,0}\int_{\bbR}\bbE\left[
\clL_{u}^{2}(t)\right] \delta_{\varepsilon }(u)du=\bbE\left[\clL_{0}^{2}(t)\right] \text{,}
\end{equation*}
by Cauchy--Schwartz, the definition of the $\delta $ function and continuity of the application $u\rightarrow \bbE[\clL_{u}^{2}(t)] $. We have thus shown that $\bbE[\clL_{\varepsilon}^{2}(t)] \rightarrow \bbE[\clL_{0}^{2}(t)],$ and the proof of Lemma~\ref{approx2} is completed.
\end{proof}

\section{Chaotic decomposition}\label{app-chaos}

We need the following standard technical result, adapted from the nodal case \cite{MPRW16} to any threshold $u\in \bbR$.

\begin{lemm}\label{indicator function}
The following decomposition holds in $L^2(\Omega)$
\begin{equation*}
\frac{1}{2\varepsilon}{1}_{[u-\varepsilon,u+\varepsilon]}(Z)=\sum_{l=0}^{+\infty }\frac{1}{l!}\beta_l^{\varepsilon}(u)\,H_{l}(Z),
\end{equation*}
where $Z\sim \clN(0,1)$, and for $l \geq 1$
\[
\beta_l^{\varepsilon}(u)=-\frac{1}{2\varepsilon}\left (
\phi \left (u+\varepsilon \right) H_{l-1} \left (u+\varepsilon \right)-
\phi \left (u-\varepsilon \right) H_{l-1} \left (u-\varepsilon \right) \right),
\]
while for $l=0$
\[
\beta_0^{\varepsilon} = \frac{1}{2\varepsilon} \int_{u-\varepsilon}^{u+\varepsilon} \phi(t)\,dt.
\]
Moreover, as $\varepsilon\to 0$,
\[
\beta_{l}^\varepsilon(u) \to \beta_l(u),
\]
where $\beta_l(u)$ coincides with~\eqref{e:beta} for every $l$.
\end{lemm}

We are now ready to establish the chaotic decomposition of the average boundary length.

\begin{proof}[Proof of Proposition~\ref{teoexpS}]
For fixed $x\in \bbS^2$, $t\in \bbR$, the projection of the random variable
\[
\frac{1}{2\varepsilon} 1_{[u-\epsilon, u+\epsilon]}(Z(x,t)) \left\|\widetilde \nabla Z(x,t)\right\|
\]
onto the chaos $C_{q}$, for $q\geq 0$, equals
\[
\sum_{m=0}^{q}\sum_{k=0}^{m} \frac{\alpha _{k,m-k}\beta^\epsilon _{q-m}(u)}{(k)!(m-k)!(q-m)!} H_{q-m}(Z(x,t))H_{k}\left(\widetilde\partial_{1;x} Z(x,t)\right)H_{m-k}\left(\widetilde{\partial}_{2;x} Z(x,t)\right),
\]
where $\{\beta_l^\epsilon(u)\}_l$ is the collection of chaotic coefficients found in Lemma~\ref{indicator function}. Since $\int_{0,T]}\int_{\bbS^2}\,dxdt<\infty$, standard arguments based on Jensen's inequality and dominated convergence yield that $\clC_T^\epsilon(u)[0] = 0$ while for $q\ge 1$
\begin{multline*}
\clC_T^{\varepsilon}(u)[q]= \sum_{m=0}^{q}\sum_{k=0}^{m} \frac{\alpha _{k,m-k}\beta^\epsilon _{q-m}(u)}{ (k)!(m-k)!(q-m)!} \\
\int_{0,T]}\int_{\bbS^2} H_{q-m}(Z(x,t))H_{k}\left(\widetilde\partial_{1;x} Z(x,t)\right)H_{m-k}\left(\widetilde{\partial} _{2;x} Z(x,t)\right)\,dx dt
\end{multline*}
in $L^2(\Omega)$. In view of Lemma~\ref{approx2} and Lemma~\ref{indicator function}, the random variable $\clC_T(u)$ being in the Wiener chaos, one has that for every $q$, as $\varepsilon \to 0$, $\clC_T^\varepsilon(u)[q]$ necessarily converge to the $q^{\rm th}$ chaotic component of $\clC_T(u)$, that is, $\clC_T(u)[q]$ as in~\eqref{e:ppS}, still in $L^2(\Omega)$.
\end{proof}




\begin{thebibliography}{99}
\bibitem{adlertaylor} Adler, Robert J.; Taylor, Jonathan E. Random Fields and Geometry, Springer Monographs in Mathematics, Springer, 2007 | MR | Zbl
\bibitem{AW:09} Azaïs, Jean-Marc; Wschebor, Mario Level Sets and Extrema of Random Processes and Fields, John Wiley \& Sons, 2009 | MR | DOI | Zbl
\bibitem{Ber02} Berry, M. V. Statistics of nodal lines and points in chaotic quantum billiards: perimeter corrections, fluctuations, curvature, J. Phys. A. Math. Gen., Volume 35 (2002) no. 13, pp. 3025-3038 | MR | DOI | Zbl
\bibitem{BP17} Berg, C. From Schoenberg coefficients to Schoenberg functions, Constr. Approx., Volume 5 (2017) no. 2, pp. 217-241 | DOI | Zbl
\bibitem{BM} Breuer, Peter; Major, Peter Central limit theorems for non-linear functionals of Gaussian fields, J. Multivariate Anal., Volume 13 (1983), pp. 425-441 | MR | DOI | Zbl
\bibitem{Ch:17} Christakos, G. Spatiotemporal Random Fields: Theory and Applications, Theory and applications, Elsevier, 2017
\bibitem{CM2018} Cammarota, Valentina; Marinucci, Domenico A quantitative central limit theorem for the Euler–Poincaré characteristic of random spherical eigenfunctions, Ann. Probab., Volume 46 (2018), pp. 3188-3288 | MR | DOI | Zbl
\bibitem{CM20} Cammarota, Valentina; Marinucci, Domenico A reduction principle for the critical values of random spherical harmonics, Stochastic Processes Appl., Volume 130 (2020) no. 4, pp. 2433-2470 | MR | DOI | Zbl
\bibitem{DM79} Dobrushin, Roland L.; Major, Peter Non-Central limit theorems for non-linear functionals of Gaussian fields, Z. Wahrscheinlichkeitstheor. Verw. Geb., Volume 50 (1979), pp. 27-52 | MR | DOI | Zbl
\bibitem{DNPR:16} Dalmao, Federico; Nourdin, Ivan; Peccati, Giovanni; Rossi, Maurizia Phase Singularities in Complex Arithmetic Random Waves, Electron. J. Probab., Volume 24 (2019), 71, 45 pages | MR | DOI | Zbl
\bibitem{EL2016} Estrade, Anne; León, José R. A central limit theorem for the Euler characteristic of a Gaussian excursion set, Ann. Probab., Volume 44 (2016) no. 6, pp. 3849-3878 | MR | Zbl | DOI
\bibitem{KL01} Kratz, Marie F.; León, José R. Central limit theorems for level functionals of stationary Gaussian processes and fields, J. Theor. Probab., Volume 14 (2001) no. 3, pp. 639-672 | MR | DOI | Zbl
\bibitem{Leo88} Leonenko, Nikolai N. On the accuracy of the normal approximation of functionals of strongly correlated Gaussian random fields, Mat. Zametki, Volume 43 (1988) no. 2, pp. 283-299 | MR | DOI | Zbl
\bibitem{LTT18} Leonenko, Nikolai N. Estimation of the covariance function of Gaussian isotropic random fields on spheres, related Rosenblatt-type distributions and the cosmic variance problem, Electron. J. Stat., Volume 12 (2018) no. 2, pp. 3114-3146 | DOI | Zbl
\bibitem{LO13} Leonenko, Nikolai N.; Olenko, Andriy Tauberian and Abelian theorems for long-range dependent random fields, Methodol. Comput. Appl. Probab., Volume 15 (2013) no. 4, pp. 715-742 | MR | DOI | Zbl
\bibitem{LP11} Lavancier, Frédéric; Philippe, Anne Some convergence results on quadratic forms for random fields and application to empirical covariances, Probab. Theory Relat. Fields, Volume 149 (2011) no. 3-4, pp. 493-514 | MR | DOI | Zbl
\bibitem{LRT17} Leonenko, Nikolai N.; Ruiz-Medina, María D.; Taqqu, Murad S. Rosenblatt distribution subordinated to Gaussian random fields with long-range dependence, Stochastic Anal. Appl., Volume 35 (2017) no. 1, pp. 144-177 | MR | DOI | Zbl
\bibitem{M2022} Marinucci, Domenico Some Recent Developments on the Geometry of Random Spherical Eigenfunctions, European Congress of Mathematics, EMS Press, Berlin, 2023, pp. 337-365 | DOI | MR | Zbl
\bibitem{MM18} Ma, Chunsheng; Malyarenko, Anatoliy Time-varying isotropic vector random fields on compact two-point homogeneous spaces, J. Theor. Probab., Volume 33 (2020) no. 1, pp. 319-339 | MR | DOI | Zbl
\bibitem{MaPeCUP} Marinucci, Domenico; Peccati, Giovanni Random Fields on the Sphere: Representations, Limit Theorems and Cosmological Applications, London Mathematical Society Lecture Note Series, 389, Cambridge University Press, 2011 | MR | DOI | Zbl
\bibitem{MPRW16} Marinucci, Domenico; Peccati, Giovanni; Rossi, Maurizia; Wigman, Igor Non-universality of nodal length distribution for arithmetic random waves, Geom. Funct. Anal., Volume 26 (2016) no. 3, pp. 926-960 | MR | DOI | Zbl
\bibitem{MRV:20} Marinucci, Domenico; Rossi, Maurizia; Vidotto, Anna Non-Universal Fluctuations of the Empirical Measure for Isotropic Stationary Fields on $\mathbb{S}^{2}\times \mathbb{R}$, Ann. Appl. Probab., Volume 31 (2021) no. 5, pp. 2311-2349 | MR | DOI | Zbl
\bibitem{MRW} Marinucci, Domenico; Rossi, Maurizia; Wigman, Igor The asymptotic equivalence of the sample trispectrum and the nodal length for random spherical harmonics, Ann. Inst. Henri Poincaré, Probab. Stat., Volume 56 (2020) no. 1, pp. 374-390 | MR | DOI | Zbl
\bibitem{noupebook} Nourdin, Ivan; Peccati, Giovanni Normal Approximations Using Malliavin Calculus: from Stein’s Method to Universality, Cambridge Tracts in Mathematics, 192, Cambridge University Press, 2012 | MR | DOI | Zbl
\bibitem{NPR19} Nourdin, Ivan; Peccati, Giovanni; Rossi, Maurizia Nodal statistics of planar random waves, Commun. Math. Phys., Volume 369 (2019) no. 1, pp. 99-151 | MR | DOI | Zbl
\bibitem{PV20} Peccati, Giovanni; Vidotto, Anna Gaussian random measures generated by Berry’s nodal sets, J. Stat. Phys., Volume 178 (2020) no. 4, pp. 996-1027 | MR | DOI | Zbl
\bibitem{Rossi2018} Rossi, Maurizia Random nodal lengths and Wiener chaos, Probabilistic Methods in Geometry, Topology and Spectral Theory (Contemporary Mathematics), Volume 739, American Mathematical Society; Centre de Recherches Mathématiques (CRM), 2019, pp. 155-169 | DOI | Zbl
\bibitem{RW08} Rudnick, Zeév; Wigman, Igor On the volume of nodal sets for eigenfunctions of the Laplacian on the torus, Ann. Henri Poincaré, Volume 9 (2008) no. 1, pp. 109-130 | MR | DOI | Zbl
\bibitem{Sze75} Szegö, Gabor Orthogonal polynomials, Colloquium Publications, 23, American Mathematical Society, 1975 | MR | Zbl
\bibitem{Ta:75} Taqqu, Murad S. Weak convergence to fractional Brownian motion and to the Rosenblatt process, Z. Wahrscheinlichkeitstheor. Verw. Geb., Volume 31 (1975), pp. 287-302 | MR | DOI | Zbl
\bibitem{Ta:79} Taqqu, Murad S. Convergence of Integrated Processes of Arbitrary Hermite Rank, Z. Wahrscheinlichkeitstheor. Verw. Geb., Volume 50 (1979), pp. 53-83 | MR | Zbl | DOI
\bibitem{VT13} Veillette, Mark S.; Taqqu, Murad S. Properties and numerical evaluation of the Rosenblatt distribution, Bernoulli, Volume 19 (2013) no. 3, pp. 982-1005 | MR | DOI | Zbl
\bibitem{Wig} Wigman, Igor Fluctuations of the nodal length of random spherical harmonics, Commun. Math. Phys., Volume 298 (2010) no. 3, pp. 787-831 | MR | DOI | Zbl
\bibitem{Wig22} Wigman, Igor On the nodal structures of random fields: a decade of results, J. Appl. Comput. Topol. (2023) | MR | Zbl | DOI
\end{thebibliography}
\end{document}
